\documentclass[12pt]{amsart}
%\usepackage[notcite,notref]{showkeys}

\usepackage{amsmath,amsxtra,amssymb,latexsym,
epsfig,amscd,amsthm,subfigure,fancybox,float,epsfig}
\usepackage[mathscr]{eucal}
\usepackage{graphicx}
\usepackage{epsfig} % for postscript graphics files
\usepackage{epstopdf}
\usepackage{cases}
%\usepackage{subfig}
%\usepackage{titlesec}
\usepackage{multicol,color}
%\titleformat{\subsection}[runin]% runin puts it in the same paragraph
%        {\normalfont\bfseries}% formatting commands to apply to the whole heading
%        {\thesubsection}% the label and number
 %       {0.5em}% space between label/number and subsection title
 %       {}% formatting commands applied just to subsection title
 %       [.]% punctuation or other commands following subsection title
\setlength{\oddsidemargin}{-0.08in}
\setlength{\evensidemargin}{-0.08in}
\setlength{\textheight}{9.0in}
\setlength{\textwidth}{6.5in}
\setlength{\topmargin}{-0.5in}
\def\para#1{\vskip .4\baselineskip\noindent{\bf #1}}
\usepackage{placeins}

\usepackage{hyperref}
\usepackage{epstopdf}
\usepackage{cases}
\usepackage{color}
%\usepackage{natbib}
\usepackage{hyperref}
\hypersetup{
    colorlinks=true,
    linkcolor=blue,
    filecolor=magenta,
    urlcolor=cyan,
    citecolor=blue,
}

%\usepackage[notcite,notref]{showkeys}
%

%\usepackage[active]{srcltx} %SRC Specials for DVI Searching
% Over-full v-boxes on even pages are due to the \v{c} in author's name
%\vfuzz2pt % Don't report over-full v-boxes if over-edge is small
%\topmargin -1cm


\newcommand{\blue}[1]{{\leavevmode\color{blue}#1}}
\newcommand{\red}[1]{{\leavevmode\color{red}#1}}
\newcommand{\cyan}[1]{{\leavevmode\color{cyan}#1}}

% THEOREM Environments --------------------------------------------------
\newtheorem{thm}{Theorem}[section]
\newtheorem {asp}{Assumption}[section]
\newtheorem{lm}{Lemma}[section]
%\newtheorem{cor}{Corollary}
\newtheorem{rmk}{Remark}[section]
\newtheorem{pron}{Proposition}[section]
\newtheorem{deff}{Definition}[section]
\newtheorem{clm}{Claim}[section]
%\newtheorem{thm}{Theorem}[section]
\newtheorem{cor}{Corollary}[section]
\newtheorem{lem}{Lemma}[section]
\newtheorem{prop}{Proposition}[section]
%\theoremstyle{definition}
\newtheorem{defn}[thm]{Definition}
\newtheorem{hyp}[thm]{Hypotheses}
\theoremstyle{remark}
\newtheorem{rem}{Remark}
\newtheorem{example}{Example}[section]
\numberwithin{equation}{section}

% MATH -----------------------------------------------------------------
\DeclareMathOperator{\RE}{Re}
 \DeclareMathOperator{\IM}{Im}
\DeclareMathOperator{\ess}{ess}
\DeclareMathOperator{\inte}{int}
\DeclareMathOperator{\suppo}{supp}
\newcommand{\eps}{\varepsilon}
\newcommand{\To}{\longrightarrow}
\newcommand{\h}{\mathcal{H}}
%\newcommand{\s}{\mathcal{S}}
\newcommand{\A}{\mathcal{A}}
\newcommand{\J}{\mathcal{J}}
\newcommand{\M}{\mathcal{M}}
\newcommand{\F}{\mathcal{F}}
\newcommand{\X}{\mathcal{X}}
\newcommand{\E}{\mathbb{E}}
\newcommand{\LL}{\mathbb{L}}
\newcommand{\N}{\mathbb{N}}
\newcommand{\PP}{\mathbb{P}}
\newcommand{\BOP}{\mathbf{B}}
\newcommand{\K}{\mathcal{K}}
\newcommand{\cB}{\mathcal{B}}
\newcommand{\BH}{\mathbf{B}(\mathcal{H})}
\newcommand{\KH}{\mathcal{K}(\mathcal{H})}
\newcommand{\R}{\mathbb{R}}
\newcommand{\Complex}{\mathbb{C}}
\newcommand{\Field}{\mathbb{F}}
\newcommand{\RPlus}{\Real^{+}}
\newcommand{\Polar}{\mathcal{P}_{\s}}
\newcommand{\Poly}{\mathcal{P}(E)}
\newcommand{\EssD}{\mathcal{D}}
\newcommand{\cC}{\mathcal{C}}
\newcommand{\Lom}{\mathcal{L}}
\newcommand{\States}{\mathcal{T}}
\newcommand{\abs}[1]{\left\vert#1\right\vert}
\newcommand{\set}[1]{\left\{#1\right\}}
\newcommand{\seq}[1]{\left<#1\right>}
\newcommand{\norm}[1]{\left\Vert#1\right\Vert}
\newcommand{\essnorm}[1]{\norm{#1}_{\ess}}
\numberwithin{equation}{section}


\newcommand{\wdt}{\widetilde}
\newcommand{\PPi}{\PP_{u,v,\pi}}
\newcommand{\Ei}{\E_{u,v,\pi}}

%\textwidth 13.5 truecm
%vsize=21.1 truecm
%\parskip 1pt

\newcommand{\op}{{\cal L}}
\newcommand{\lbar}{\overline}
\newcommand{\bed}{\begin{displaymath}}
\newcommand{\eed}{\end{displaymath}}
\newcommand{\bea}{\bed\begin{array}{rl}}
\newcommand{\eea}{\end{array}\eed}
\newcommand{\ad}{&\!\!\!\disp}
\newcommand{\aad}{&\disp}
\newcommand{\barray}{\begin{array}{ll}}
\newcommand{\earray}{\end{array}}
\newcommand{\diag}{{\rm diag}}
\def\disp{\displaystyle}
\newcommand{\rr}{{\Bbb R}}
\newcommand{\1}{\boldsymbol{1}}

\newcommand{\bmu}{\boldsymbol{\mu}}
\newcommand{\bdelta}{\boldsymbol{\delta}}


 \font\inh=cmr10 \font\cto=cmr10
\font\tit=cmbx12 \font\tmd=cmbx10 \font\ab=cmti9 \font\cn=cmr9
\def\ker{\text{Ker}}
\def\bar{\overline}
\def\hat{\widehat}
\def\a.s{\text{\;a.s.\;}}
\def\supp{\text{supp\,}}
%\pagestyle{headings}


\title{%A
Stochastic Epidemic SIR Models with Hidden States}



\author[N. Du]{Nguyen Du }
\address{Department of Mathematics-Mechanics-Informatics \\
Vietnam National University of Science\\
34 Nguyen Trai\\
 Thanh Xuan, Ha Noi\\
Vietnam}
\email{nhdu@viasm.edu.vn}

\author[A. Hening]{Alexandru Hening }
\address{Department of Mathematics\\
Texas A\&M University\\
Mailstop 3368\\
College Station, TX 77843-3368\\
United States
}
\email{ahening@tamu.edu}

\author[N. Nguyen]{ Nhu Nguyen }
\address{Department of Mathematics \\
University of Connecticut\\
341 Mansfield Road U1009\\
Storrs, Connecticut 06269-1009\\
United States. {\rm The research of this author was supported in part by the National Science Foundation.}}
\email{nguyen.nhu@uconn.edu }


\author[G. Yin]{George Yin}
\address{Department of Mathematics \\
University of Connecticut\\
341 Mansfield Road U1009\\
Storrs, Connecticut 06269-1009\\
United States. {\rm The research of this author was supported in part by the National Science Foundation.}}
\email{gyin@uconn.edu}


\keywords {SIR model; Extinction;  Permanence; Stationary Distribution; Ergodicity}
\subjclass[2010]{34C12, 60H10, 92D25}

\begin{document}
\begin{abstract}
%We
This paper focuses on and
analyzes  realistic SIR models that take  stochasticity
%which takes
into account. The proposed systems
%Our results
are applicable to
%valid  for
most incidence rates that are used in the literature including the bilinear incidence rate, the Beddington-DeAngelis incidence rate, and a Holling type II functional response.
 Given that many diseases can lead to asymptomatic infections, we look at a system of stochastic differential equations that also includes a class of hidden state individuals, for which the infection status is unknown. We assume that the direct observation of the percentage of hidden state individuals that are infected, $\alpha(t)$, is not given and only
 %we  can only observe
 a noise-corrupted observation
 %noisy
 process is available.
 %of the type $dy(t)=g(\alpha(t))dt+dW(t)$.
 %We are able to use
 Using
 the nonlinear filtering techniques
 %model
 in conjunction with an invasion type analysis (or analysis using Lyapunov exponents from the dynamical system point of view), this paper
 %in order
 %to
 proves that the long-term behavior of the disease is governed by a threshold $\lambda\in \R$
  that depends on the model parameters. It turns out that if $\lambda<0$ the number $I(t)$ of infected individuals converges to zero exponentially fast, or the extinction happens. In contrast,
  if $\lambda>0$, the infection is endemic and the system %converges to its unique invariant probability measure
  is permanent.
  We showcase our results by applying them in specific illuminating examples.
  %settings.
Numerical simulations are also given to illustrate our results.
\end{abstract}

%\setlength{\baselineskip}{0.28in}
\maketitle


\section{Introduction}\label{sec:int}


The SIR epidemic models introduced first by \cite{KER1,KER2} look at the dynamics of susceptible, infected, and recovered individuals, whose densities at the time $t$ are denoted by $S(t)$, $I(t)$, and $R(t)$, respectively. In the absence of random effects,
%If random effects are not considered,
the dynamics  are
%in general
described
by the following system of differential equations
\begin{equation}\label{1.0}
\begin{cases}
dS(t)=\big[a_1-\mu_SS(t)-F(S(t),I(t))\big]dt, \\
dI(t)=\big[-(\mu_I+r)I(t) + F(S(t),I(t)\big]dt,\\
dR(t)=\big[-\mu_RR(t)+rI(t)\big]dt.
\end{cases}
\end{equation}
Here $a_1>0$ is the recruitment rate of the population, $\mu_S,\mu_I,\mu_R>0$ are the death rates of the susceptible, infected, and recovered individuals, $r>0$ is the recovery rate of the infected individuals and $F(S(t),I(t))$ is the incidence rate.
The dynamics of recovered individuals have no effect on that of the disease transmission.
% dynamics.
As a result, it is the usual practice not to consider the recovered individuals as part of the problem formulation. We adopt this practice
%-- we will do this
throughout this paper.
 Various types of incidence rates $F(S,I)$ have been considered in the literature, for example,
the Holling type II functional response  $F(S,I)=\frac{\beta SI}{m_1+S}$ (see \cite{HUO}), the bilinear functional response $F(S,I)=\beta SI$ (see \cite{DANG,ZHANG}), the nonlinear functional response  $F(S,I)=\frac {\beta SI^l}{1+m_2I^h}$ (see \cite{RUAN,YANG}),
and the Beddington-DeAngelis functional response $F(S,I)=\frac{\beta SI}{1+m_1S+m_2I}$  (see \cite {NHU1, NHU2}).
%\end{itemize}
%For the deterministic SIR models with these incidence rates, the researchers have found the reproduction number $R_0$ which has the property: if $R_0<1$ then the disease free equilibrium point is locally asymptotically stale; in case $R_0>1$ we see that the disease point is unstable and there is a steady state, which is locally asymptotically stable.
%We can find those results in \cite{Kad1, Kad2} for the deterministic SIR model with the Beddington-DeAngelis incidence rate $\dfrac{\beta SI}{1+m_1S+m_2I}$ or in \cite{PZHANG} for the  bilinear incidence $\beta SI$, or in \cite{CAI,MA,SHAN,SUN,WANG} for some other forms of the incidence rate.

It is by now widely known that in order
%If one wants
to have
a realistic model, one cannot ignore random environmental fluctuations (temperature, the climate, the water resources, etc.). In this paper we consider stochastic epidemic models of the form
%By these random effects, the death rates and the incidence rate are often perturbed by white noises. Many authors have considered the stochastic SIR models when the natural death rates are affected by white noises, i.e., $\mu_S \hookrightarrow \mu_S + \sigma_1\dot B_1(t), \mu_I+r\hookrightarrow \mu_I+r+\sigma_2 \dot B_2(t)$, where $B_i(t),i=1,2$ are Brownian motions and the stochastic equation in general has the form (see e.g., \cite{DANG,NHU2,YANG})
\begin{equation}\label{e:unhid}
\begin{cases}
dS(t)=\big[a_1-b_1S(t)-I(t)f(S(t),I(t))\big]dt + \sigma_1 S(t) dB_1(t),\\
dI(t)=\big[-b_2I(t) + I(t)f(S(t),I(t))\big]dt + \sigma_2I(t) dB_2(t),\\
\end{cases}
\end{equation}
where $B_1(t)$ and $B_2(t)$ are independent Brownian motions and $\sigma_1,\sigma_2\neq 0$ are the noise intensities (standard deviations).
%variances of the noise.
Moreover, in the above we have rewritten the coefficients: $b_1:=\mu_S, b_2:=\mu_I+r$ and $F(S,I)=If(S,I)$ (compare this to \eqref{1.0}). This system has been analyzed in a
%very
general setting in \cite{DN-JDE}.

It is well-known that there are diseases for which certain infected individuals
%while infected,
are asymptomatic.
Covid-19 is one such example -- there have been many reports of infections where the infected exhibit no symptoms. We intend to capture this type of behavior in our model.
In order to do this, we assume that
the group of infected individuals that has incidence rate $I(t)f(S(t),I(t))$ (the rate that describes how the disease spreads from infected groups to susceptible groups) in the classical setting, now contains 2 sub-groups. The first group contains individuals who have been confirmed to be infected and
%has
with the incidence rate $I(t)f_1(S(t),I(t))$ (we will still denote this incidence rate by $I(t)f(S(t),I(t))$
%later
 for
 %the
  notational simplicity).
%   of notation).
The second group has incidence rate $I(t)h(S(t),I(t))$ and contains  people whose
%there is fourth population class, with density at time $t$ denoted by $H(t)$.
%If an individual is in this new class, their
infection status is unknown or hidden.
Let $\alpha(t)$ be a Markov process
%chain
taking values in $\M=[0,1]$. We suppose that $\alpha(t)$
%will
represents the percentage of
%the rate of
individuals in
the hidden-status class
%individuals
that are infected at time $t\geq0$
%We model the dynamics of the density $H(t)$ of this class by
%$$
%dH(t)=\bar h(\alpha(t),S(t),I(t))dt
%$$
%where $\bar h(\alpha,S,I)=Ih(\alpha,S,I)$.
%and having generator $Q=(q_{ij})_{i,j\in\M}$
%However, we can not observe directly this class. In other word, we do not know that is infected or not.
%On the other hand, although, we can not observe $\alpha(t)$, we can observe partially by noise via a observation process $y(t)$.
%Assume we known $y(t)$ with
%$$
%dy(t)=g(\alpha(t))dt+dW(t).
%$$
and that only noise-corrupted observations of
%the direct
%observation of the process
$\alpha(t)$
are
%not
available. More specifically, one can only observe $\alpha(t)$ with additive white noise.
%, and that data concerning $\alpha(t)$ is only observed via noise through
%the process $y(t)$, which will be given later in \eqref{eq-oby}.

It is natural to assume that the hidden status of potentially infected individuals affects the spread of the disease. As a result,
we let
 the functions $f$ and $h$
%will
depend on $\alpha(t)$.
%As a result, including
With
the hidden state dynamics in \eqref{e:unhid}, we obtain
\begin{equation}\label{eq-main-1}
\begin{cases}
dS(t)=\Big[a_1-b_1S(t)-I(t)f(\alpha(t),S(t),I(t))-\alpha(t)I(t)h(\alpha(t),S(t),I(t))\Big]dt + \sigma_1 S(t) dB_1(t),\\[1ex]
dI(t)=\Big[-b_2I(t) + I(t)f(\alpha(t),S(t),I(t))+\alpha(t)I(t)h(\alpha(t),S(t),I(t))\Big]dt + \sigma_2I(t) dB_2(t).
\end{cases}
\end{equation}

\begin{rem}
One can understand the dynamics by looking at individuals from group $S$, in which
%Infected individuals infect
susceptible individuals are infected at the rate $If(S,I)$.
%Someone can be potentially infected, i.e. a member of the hidden group.
We assume that $\alpha(t)$ percent of the rate of the potentially infected individuals are actually infected. Then we can say that members of the hidden group infect susceptible individuals at the rate $\alpha I h(S,I)$.
\end{rem}
\begin{rem}
%	Note that from a mathematical point of view,
We could combine $I(t)f(\alpha(t),S(t),I(t))+\alpha(t) I(t) h(\alpha(t),S(t),I(t))$ to
produce  a new function. However, we choose the current setup
%want
 to make
the formulation and motivation clear. Moreover, this will also be more convenient for later discussions.
\end{rem}

Our results can be summarized as follows.
%Note first that since
Because the infection status of certain individuals is hidden,
%since
and
$\alpha(t)$ is not directly available,
%observable,
the dynamics of \eqref{eq-main-1} are difficult to study.
%In order to bypass this
To overcome the
difficulty, we apply the nonlinear filtering techniques by considering the conditional distribution of the process $\alpha(t)$ given the observations. This enables us to
%formulate in
replace the hidden Markov process $\alpha(t)$ in
\eqref{eq-main}
by
the corresponding conditional distribution.
%corresponding ``observable" stochastic differential equation version of
%\eqref{eq-main-1}.
We start by studying the well-posedness of the equation under consideration together with the positivity of solutions, the Markov-Feller property, and some moment estimates. Next, we study the longtime behavior of the system.
Under the assumption for ergodicity of nonlinear filtering  \cite{Han09,Kun71} and using ideas from dynamical systems, by considering the boundary equation and growth rate (see e.g., \cite{DANG,DN-JDE,HN18}), we are able to prove that
%the following result:
there is a threshold $\lambda$ such that if $\lambda<0$, the number of the infected individuals $I(t)$ tends to zero
%at an
exponentially fast
%rate
and if $\lambda>0$, all invariant probability measures of the system concentrate on $\mathbb{R}_+^{2,\circ}:=(0,\infty)^2$, and then the systems is permanent.
We show that the threshold $\lambda$ also characterizes
%is in fact the one characterizing
the permanence and extinction of the system \eqref{eq-main-1}.
%In other words, the observable equation \eqref{eq-main} that uses the nonlinear filter, %preserves the characterization of \eqref{eq-main-1}.
 We also study the case when the process $\alpha(t)$ is a hidden Markov chain
 taking values in a finite set.
 %We end by looking some
 Next, we demonstrate our results using
 simple examples and numerical simulations.
 % which help to illustrate our theoretical results.
%Some discussion about relaxing conditions and future research are also provided.


The rest of the paper is organized as follows. We give the mathematical formulation of our problem in Section \ref{secpre}.
Section \ref{sec:thres} introduces the threshold $\lambda$. The sign of $\lambda$ will be used to characterize the longtime behavior of the underlying system.
Section \ref{sec:lon} is devoted to the characterization of the longtime dynamics of our system. Section \ref{sec:disc} offers
%various
some interpretations and implications of our results.
Finally, Section \ref{sec:exp} provides some simple examples and simulations to illustrate our theoretical results.

\section {%Filtering
Problem
%and
Formulation}\label{secpre}

Throughout this paper we use
$\R_+:=[0,\infty)$, $\R^\circ_+:=(0,\infty)$,
$\R_+^2:=[0,\infty) \times [0,\infty)$, and $\R^{2,\circ}_+:=(0,\infty)\times (0,\infty)$.
Let $(\Omega,\F,\{\F_t\}_{t\geq0},\PP)$ be a complete probability space with filtration $\{\F_t\}_{t\geq0}$ satisfying the usual conditions, and $B_1(t)$, $B_2(t)$, and $W(t)$ be mutually independent standard Brownian motions. The process
$\alpha(t)$ (termed a signal process) is assumed to be an adapted stochastic process taking values in $[0,1]$ that is independent of $B_1(t)$, $B_2(t)$, and $W(t)$.
Moreover, $\M$ will denote the space of probability measures on $([0,1],\mathcal B ([0,1]))$ endowed with the weak topology,
%and
%$\mathcal P(\M)$ is the space of probability measure on $\M$,
and $C[0,1]$ the spaces of all real-valued continuous functions on $[0,1]$.
For any function $l\in C[0,1]$ and $\mu\in\M$, set
$$
\mu(l):=\int_0^1 l(x)\mu(dx).
$$

As discussed above, we consider the setting where
%direct  observations of
the precise values
$\alpha(t)$ are not available and only noisy observations are available.
%possible.
The observation process $y(t)$ of the signal process $\alpha(t)$
%under noise
is given by
\begin{equation}\label{eq-oby}
dy(t)=g(\alpha(t))dt+dW(t),\quad y(0)=0,
\end{equation}
where $g:[0,1]\to\R$ is a continuous function.
Let
$
\F_t^y:=\sigma\{y(s):0\leq s\leq t\}\bigvee\sigma(\alpha(0)),$
where $\bigvee$ denotes the smallest $\sigma$-algebra generated by the union of some $\sigma$-algebras.
Let $\Pi_t(\cdot)\in \M$ be the conditional distribution of the signal process $\alpha(t)$ given the observation $y(t)$ and the initial data, i.e.,
$$
\Pi_t(A)=\PP[\alpha(t)\in A|\F_t^y], A\in\mathcal B([0,1]).
$$
Such $\{\Pi_t(\cdot)\}$ is called nonlinear filtering.

% Moreover, to simplify notation, we will often work with some fixed initial distribution of $\alpha(0)$ and suppress the dependence of $\Pi_t(\cdot)$ on the initial distribution of $\alpha(0)$.

The field of nonlinear filtering has a long history.
The main idea stems from replacing  the unknown state by its conditional distributions. The earliest result was the well-known Kushner's equation
\cite{Kushner64}. Subsequently, the Duncan-Mortensen-Zakai equation came into being
\cite{Duncan,Mortensen,Zakai}. In this paper, we make use of the version of filtering developed by
Fujisaki-Kallianpur-Kunita \cite{FFK72}. We will also make use of the Wonham filter for hidden Markov chains, which is one of the handful finite-dimensional filters in existence \cite{Won65}.

To proceed,
we detail the results of
Fujisaki-Kallianpur-Kunita \cite{FFK72} (see also \cite{KS69}), which involve
%obtained a stochastic
a differential equation for the nonlinear filtering $\Pi_t(\cdot)$, next.
% in a general setting.
Define
$$
\beta(t)=y(t)-\int_0^t \Pi_s(g)ds,
$$
and note that the process $\beta(t)$ is a one-dimensional Wiener process, see e.g., \cite[Theorem 7.2]{LS01}.
Moreover, $
\sigma\{\beta(t_2)-\beta(t_2):t_2\geq t_1\geq t\}$ and $\F_t^y$ are independent for all $t\geq 0.$
If the signal process $\alpha(t)$ is a Markov process with infinitesimal generator $A$, then $\Pi_t(\cdot)$ is the solution of
\begin{equation}\label{eq-pi}
\Pi_t(l)=\Pi_0(l)+\int_0^t\Pi_s(Al)ds+\int_0^t \left( \Pi_s(lg)-\Pi_s(l)\Pi_s(g)\right)d\beta(s),\quad\forall l\in\mathcal D(A).
\end{equation}
The interested reader is referred to the detailed analysis given in \cite{FFK72,KS69}.

We will not make use of \eqref{eq-pi} often in our analysis, except for establishing some preliminary properties.
%This is because
The stochastic differential equation for $\Pi_t(\cdot)$ is
%very complicated
rather complex and is not the main concern of the current paper.
% to analyze.
%More importantly, such an explicit equation is also not necessary in our analysis.
As will be seen in the next section, we only need
 %require to
to establish the related ergodicity. Thus,
for us, it suffices to consider
%is enough to see
$\Pi_t(\cdot)$ as a stochastic process taking values in $\M$.
In addition,
%We will  consider the
we use
continuous measurable modification of $\Pi_t(\cdot)$;
% and note that
such a modification always exists
%(see
\cite{Kun71}.

Under the premise that
%We can only observe
one only observes a noisy version of
$\alpha(t)$,
% and
we proceed to study  system \eqref{eq-main-1} by using
%. We will use
the nonlinear filter $\{\Pi_t(\cdot)\}$
%to study \eqref{eq-main-1}
with given the information of the observation process $y(t)$. More precisely, we consider the system
\begin{equation}\label{eq-main}
\begin{cases}
dS(t)=\Big[a_1-b_1S(t)-\int_0^1I(t)f(x,S(t),I(t))\Pi_t(dx)
\\
\hspace{2cm}
-I(t)\int_0^1xh(x,S(t),I(t))\Pi_t(dx)\Big]dt + \sigma_1 S(t) dB_1(t),\\[1ex]
dI(t)=\Big[-b_2I(t) + \int_0^1I(t)f(x,S(t),I(t))\Pi_t(dx)
\\
\hspace{2cm}
+I(t)\int_0^1xh(x,S(t),I(t))\Pi_t(dx)\Big]dt + \sigma_2I(t) dB_2(t),
\end{cases}
\end{equation}
where
$$
\begin{aligned}
&\Pi_t(A)=\PP[\alpha(t)\in A|\F_t^y],\forall A\in \mathcal B([0,1]),\\
&dy(t)=g(\alpha(t))dt+dW(t), \quad y(0)=0.
\end{aligned}
$$
%We will
Denote by $\PPi$ and $\Ei$ the probability and expectation corresponding to the initial values $S(0)=u$, $I(0)=v$,  $\Pi_0=\pi$, and the distribution of $\alpha(0)$, respectively.
We next make some assumptions that will be used throughout this paper.
% the solution to \eqref{eq-main} with initial values $S(0)=u, I(0)=v$.

%As said before, although $\Pi_t(\cdot)$ depends on $\alpha(0)$ and then so are $S(t)$ and $I(t)$, we will often fix $\alpha(0)$ and suppress this dependence.

%We next make some assumptions that will be used throughout this paper.
%Throughout of this paper, we assume that the incidence rate
%and  the diffusion term
%satisfy the following conditions.
\begin{asp}\label{asp-1}
	The following conditions hold:
	\begin{itemize}
		\item The function $f:[0,1]\times \R^2_+\to\R_+$ is nonnegative, $f(x,0,i)=0, \forall x\in[0,1]$, $i\geq 0$. Furthermore, $f$ is Lipschitz continuous,
		%in $s,i$ uniformly in $x$,
		i.e., there exists a positive constant $L_1$ such that
		 for all $ x_1,x_2\in[0,1], s_1,s_2,i_1,i_2\geq 0$
		\begin{equation*}
		\begin{aligned}\abs{f(x_1,s_1,i_1)-f(x_2,s_2,i_2)}\leq L_1(|x_1-x_2|+\abs{s_1-s_2}+\abs{i_1-i_2}).
		\end{aligned}
		\end{equation*}
		\item The function $h:[0,1]\times\R_+^2\to\R_+$ satisfies $h(x,0,i)=0$, and is Lipschitz continuous with Lipschitz constant $L_2$, i.e., for all $s_1,s_2,i_1,i_2\geq 0, x_1,x_2\in[0,1]$,
				\begin{equation*}
		\begin{aligned}\abs{h(x_1,s_1,i_1)-h(x_2,s_2,i_2)}\leq L_2(|x_1-x_2|+\abs{s_1-s_2}+\abs{i_1-i_2}).
		\end{aligned}
		\end{equation*}
%		and there is $\kappa>0$ such that
%		\begin{equation}\label{eq-conh}
%		\inf_{(x,s,i)\in [0,1]\times\R^2_+}(b_1-\kappa)s+(b_2-\kappa)i-ih(x,s,i)\geq 0.
%		\end{equation}
	\item For each $x\in[0,1]$, functions $h(x,\cdot,0)$ and $f(x,\cdot,0)$ are non-decreasing.
	\end{itemize}
\end{asp}

\begin{rem}
We note that almost all of the incidence rate functions used in the literature (such as the bilinear incidence rate, the Beddington-DeAngelis incidence rate, the Holling type II functional response, etc.) satisfy these conditions. \textit{Recall that the incidence rate in our setting is $If(S,I)$ rather than $f(S,I)$}.
%Moreover, the boundedness of $h$ is not restrictive. Note that the growth rate of hidden class is $\bar h(\alpha,I)=Ih(\alpha,I)$.
%Condition \eqref{eq-conh} is motivated by applications. We see that $``b_1s+b_2i"$ is the ``moving out" rate (by death or moving to hidden status class) of susceptible and infected classes. It is clear that this quantity is strictly larger than the ``moving in" rate of the hidden class.
%\textit{Note that the growth rate of the hidden class is $\bar h(\alpha,S,I)=Ih(\alpha,S,I)$ rather than $h(\alpha,S,I)$}.

The third condition is imposed because the incidence rate and the growth rate of the hidden class should increase when $S(t)$ and $I(t)$ increase. Since we have rewritten these rates as $If(S,I)$ and $Ih(\alpha,S,I)$, only increasing condition on $S$ is assumed.
\end{rem}

\begin{asp}\label{asp-2}
	The signal process $\alpha(t)$ is a Markov-Feller process that has a unique invariant measure $\mu^*$ and
		$$
		\|P(t,x,\cdot)-\mu^*(\cdot)\|_{\text{TV}}=0,
		$$
		where $P(t,x,\cdot)$ is the transition probability and $\|\cdot\|_{\text{TV}}$ is the total variation norm.
\end{asp}

\begin{rem}
	This assumption is needed to guarantee the ergodicity of the nonlinear filtering as discussed in Section \ref{sec:filter} later.
	Using this, we can then define a threshold that fully characterizes the longtime behavior of the underlying system; see Section \ref{sec:lambda}.
\end{rem}

%\red{NOTE: Here we need to explain in a Remark why this assumption is required}


For $V(s,i):\R^2\to\R$, define the operator $\Lom V$ by
$$
\begin{aligned}
\Lom V[s,i,\pi]=&\dfrac{\partial V}{\partial s} \Big[a_1-b_1s-if(s,i)-i\int_0^1xh(x,s,i)\pi(dx)\Big]\\
&+\dfrac{\partial V}{\partial i} \Big[-b_2i+if(s,i)+i\int_0^1xh(x,s,i)\pi(dx)\Big]\\
&+\frac{\sigma_1^2s^2}2\dfrac{\partial^2 V}{\partial s^2}+\frac{\sigma_2^2i^2}2\dfrac{\partial^2 V}{\partial i^2}.
\end{aligned}
$$
In the above, $(s,i,\pi)\in\R^2\times \M$ represents the variable of $\Lom V$ rather than that of $V$.



\para{Discrete state space and Wonham filter.}
If the Markov process $\alpha(t)$ takes values in a finite space $\{m_1,\dots,m_{n^*}\}\subset [0,1]$ and has generator $\{q_{ik}\}_{i,k\in\{1,\dots,n^*\}}$, the formulation will be
%more
simpler and  more explicit.
We can formulate the problem as follows.
Let
$$
\begin{aligned}
%& \F_t^y:=\sigma\{y(s):0\leq s\leq t\},\\
& e_k(t):=\PP(\alpha(t)=m_k|\F_t^y)=\E [\1_{\{\alpha(t)=m_k\}}|\F_t^y], k=1,\dots,n^*,\\ &e(t)=(e_1(t),\dots,e_{n^*}(t)),\\
& \mathcal S_{n^*}:=\left\{e=(e_1,\dots,e_{n^*})\in\R^{n^*}:e_k\geq 0,\sum_{k=1}^{n^*}e_k=1\right\},\\
& g_k:=g(m_k), k=1,\dots,n^*, \quad \bar g(e):=\sum_{k=1}^{n^*} g_ke_k, e=(e_1,\dots,e_{n^*})\in\mathcal S_{n^*}.
\end{aligned}
$$
It was shown in \cite{Won65} that the posterior probability $e(t)$
satisfies the following system of stochastic differential equations
\begin{equation}
\begin{cases}
de_k(t)=\disp\!\left[\!\sum_{i=1}^{n^*} q_{ik}e_i(t)-(g_k-\bar g(e(t)))\bar g(e(t))e_k(t) \!\!\right]\!\!dt\!\!+\! (g_k-\bar g(e(t)))e_k(t)dy(t),\ k=1,\dots,n^*,\\
e_k(0)=e_k^0,\quad k=1,\dots,n^*.
\end{cases}
\end{equation}
In this case, instead of considering  system \eqref{eq-main}, one can study the following system of stochastic differential equation
\begin{equation}\label{eq-main-dis}
\begin{cases}
dS(t)=\Big[a_1-b_1S(t)-I(t)\sum_{k=1}^{n^*}f(m_k,S(t),I(t))e_k(t)
\\
\hspace{2cm}
-I(t)\sum_{k=1}^{n^*}m_kh(m_k,S(t),I(t))e_k(t)\Big]dt+ \sigma_1 S(t) dB_1(t),\\[1ex]
dI(t)=\Big[-b_2I(t) + I(t)\sum_{k=1}^{n^*}f(m_k,S(t),I(t))e_k(t)
\\
\hspace{2cm}
+I(t)\sum_{k=1}^{n^*}m_kh(m_k,S(t),I(t))e_k(t)\Big]dt + \sigma_2I(t) dB_2(t),\\
de_k(t)=\left[\sum_{i=1}^{n^*} q_{ik}e_i(t)-(g_k-\bar g(e(t)))\bar g(e(t))e_k(t) \right]dt+(g_k-\bar g(e(t)))e_k(t)dy(t),\quad  k=1,\dots,n^*.\\
%\varphi_k(0)=\varphi_k^0,\quad k=1,\dots,n^*.
\end{cases}
\end{equation}
Moreover, the process
$$
\bar W(t)=y(t)-\int_0^t \bar g(e(s))ds
$$
is a one-dimensional Brownian motion adapted to $\F_t^y$; see e.g., \cite[Theorem 7.2]{LS01}. Therefore,  system \eqref{eq-main-dis} can be rewritten as
\begin{equation}\label{eq-main-dis-1}
\begin{cases}
dS(t)=\Big[a_1-b_1S(t)-I(t)\sum_{k=1}^{n^*}f(m_k,S(t),I(t))e_k(t)
\\
\hspace{2cm}
-I(t)\sum_{k=1}^{n^*}m_kh(m_k,S(t),I(t))e_k(t)\Big]dt + \sigma_1 S(t) dB_1(t),\\[1ex]
dI(t)=\Big[-b_2I(t) + I(t)\sum_{k=1}^{n^*}f(m_k,S(t),I(t))e_k(t)
\\
\hspace{2cm}
+I(t)\sum_{k=1}^{n^*}m_kh(m_k,S(t),I(t))e_k(t)\Big]dt + \sigma_2I(t) dB_2(t),\\
de_k(t)=\sum_{i=1}^{n^*} q_{ik}e_i(t)+e_k(t)(g_k-\bar g(e(t)))d\bar W(t),\quad  k=1,\dots,n^*.\\
%\varphi_k(0)=\varphi_k^0,\quad k=1,\dots,n^*.
\end{cases}
\end{equation}
This system is easier to analyze than \eqref{eq-main}. However, in this case, we need to assume that the signal process $\alpha(t)$ representing the portion of the rate of the infection in the group of individuals with hidden infection status takes only finitely many values.
%[QQ: $H$ defined??]
This significantly limits the setting as well as the possible applications to real world problems.

One may simplify the problem further by assuming that $\alpha(t)$ takes values in $\{0,1\}$. This would mean that at any given time all individuals in the hidden status group are either susceptible or infected.

\section{Ergodicty of Nonlinear Filter and Threshold for Permanence and Extinction}\label{sec:thres}
\subsection{Ergodicity of Nonlinear Filter}\label{sec:filter}
The study of the asymptotic properties of the nonlinear filter has a long history in the literature. We briefly summarize the developments. One of the first works is Kunita's paper \cite{Kun71}.
We restate the main result (Theorem 3.3) of this reference as follows.

\begin{pron}
%{thm}
$($Kunita 1971$)$
	Assume that the signal process $\alpha(t)$ taking values in a compact separable Hausdorff space is a Markov-Feller process with semigroup $P_t$ that has a unique invariant measure $\mu^*$ and
	$$
	\limsup_{t\to\infty}\int|P_tl(x)-\mu^*(l)|\mu(dx)=0,\forall l\in C([0,1]).
	$$
	Then the process $\Pi_t(\cdot)$ is an $\M$-valued Markov-Feller process that has unique invariant measure $\Phi^*$. Moreover, $\mu^*$ is the barycenter of $\Phi^*$, i.e.,
	$$
	\mu^*(l)=\int_{\M} \nu(l)\Phi^*(d\nu), \forall l\in C[0,1].
	$$
\end{pron}
%{thm}

Unfortunately, it was pointed out in \cite{BCL04} that there was a serious gap in the proof of the main result in \cite{Kun71}. A key role in the verification of the uniqueness for the invariant measure of $\Pi_t(\cdot)$ is the following identity
\begin{equation}\label{eq-identity}
\bigcap_{t\geq 0}\F_{[0,\infty)}^y\bigvee \sigma\{\alpha(s):s\geq t\}=\F_{[0,\infty)}^y\bigvee \left(\cap_{t\geq 0}\sigma\{\alpha(s):s\geq t\}\right),
\end{equation}
where $\F_{[0,\infty)}^y:=\bigvee_{t\geq 0}\F_t^y$.
This identity is indispensable in the proof of the uniqueness of the invariant measure of nonlinear filtering; see the counterexample given in \cite{BCL04}.
Moreover, the exchange of intersection and supremum is not always permitted in general\footnote{According to Williams \cite{Wil91}, this incorrect identity ``...tripped up even Kolmogorov and Wiener"; see \cite[p. 837]{Sin89}, and \cite[pp. 91--93]{Mas66}}.
However, in Kunita's proof, this identity was not proved.
%Indeed, at the crucial point in the proof ([12], top of p. 384), Kunita implicitly takes for granted that the following exchange of intersection and superemum
On the other hand, it is important to note that all the known counterexamples are based on the degeneracy of the observation, i.e., there is no added noise.
Therefore, it was tempting to conjecture that the identity  \eqref{eq-identity} still holds provided the nondegeneracy of the observation.

In 2009, R. Handel \cite{Han09}  has
%completed
partially solved this open problem
%to solve
 in a general setting.
In fact, \cite[Theorem 4.2]{Han09} proved that  identity \eqref{eq-identity} does indeed
hold under conditions of ergodicity of signal process \cite[Assumption 3.1]{Han09} and nondegeneracy of the observation process \cite[Assumption 3.2]{Han09}, which are only mildly stronger than those in \cite{Kun71}\footnote{According to Handel \cite{Han09}, whether Kunita's condition is already sufficient to guarantee uniqueness of the
	invariant measure with barycenter $\mu^*$ remains an open problem.}.
Finally, we state the following theorem on the ergodicity of the filter $\Pi_t(\cdot)$ under our setting and our assumption.

\begin{pron}
%{thm}
	Under Assumption \ref{asp-2}, the process $\Pi_t(\cdot)$ is an $\M$-valued Markov-Feller process and has a unique invariant measure $\Phi^*$. Moreover, $\mu^*$ is the barycenter of $\Phi^*$, i.e.,
	$$
	\mu^*(l)=\int_{\M} \nu(l)\Phi^*(d\nu), \forall l\in C[0,1].
	$$
\end{pron}
%{thm}
%Moreover, a simpler result for ergodicity of Wonham filter in case of finite state space can be found in \cite[Section 3]{TW11}.

Moreover, let $\M_{\Phi^*}\subset\M$ be the support of the invariant measure $\Phi^*$ of the nonlinear filter $\Pi_t(\cdot)$. In general, one should not expect that $\M_{\Phi^*}=\M$. In fact, this does not hold even in the simple setting of the Wonham filter when the state space has only 3 states \cite[Section 4]{TW11}.

\subsection{Threshold for Permanence and Extinction}\label{sec:lambda}
%\subsection{The threshold $R$}
We next use the ergodicity of the nonlinear filter developed in the previous section in conjunction with a Lyapunov exponent analysis (sometimes called invasion type analysis in population dynamics) coming from dynamical systems \cite{DANG,DN-JDE,HN18}. This allows us to introduce a threshold $\lambda$, which characterizes the longtime behavior of system \eqref{eq-main}.

Consider the equation on boundary when the infected individuals are absent, i.e.,
\begin{equation}\label{phi}
d \varphi(t)=\big(a_1-b_1\varphi(t)\big)dt+\sigma_1 \varphi(t)dB_1(t),\;\;\;\varphi(0)=u\geq 0.
\end{equation}
%We write $\varphi_u(t)$ for the solution of the equation \eqref{phi} with the initial condition $\varphi(0)=u$.
By solving the Fokker-Planck equation,   equation \eqref{phi}  has
a unique stationary distribution $\hat\mu$ with the density  given by
\begin{equation}\label{hpp}
\dfrac{b^a}{\Gamma(a)}y^{-(a+1)}e^{-\frac{b}y},\; y>0,
\end{equation}
where $c_1=b_1+\dfrac{\sigma_1^2}{2}, a=\dfrac{2c_1}{\sigma_1^2},b=\dfrac{2a_1}{\sigma_1^2} $ and
$\Gamma(\cdot) $ is the Gamma function.
The main idea is to determine whether $I(t)$ converges to 0 or not by
looking at the Lyapunov exponent $\limsup_{t\to\infty} \dfrac{\ln I(t)}t$ when $I(t)$ is small. Using It\^{o}'s formula yields
\begin{equation}\label{Lia}
\begin{aligned}
\dfrac{\ln I(t)}{t}=&\dfrac{\ln v}{t}+ \dfrac {\sigma_2B_2(t)}t-c_2+\dfrac 1{t}\int_0^t\int_0^1f(x,S(s),I(s))\Pi_s(dx)ds\\
&+\frac 1t\int_0^t \int_0^1xh(x,S(s),I(s))\Pi_s(dx)ds,
\end{aligned}
\end{equation}
where $c_2=b_2+\dfrac{\sigma_2^2}2.$
Intuitively, $\limsup_{t\to\infty} \frac{\ln I(t)}t<0$ implies $\lim_{t\to\infty} I(t)=0$. As a result, if $I(t)$ is small then $S(t)$ is close to $\varphi(t)$ provided $S(0)=\varphi(0)$. Therefore, when $t$ is sufficiently large we have
\begin{equation*}
\begin{aligned}
\dfrac 1{t}&\int_0^t\int_0^1f(x,S(s),I(s))\Pi_s(dx)ds+\frac 1t\int_0^t \int_0^1xh(x,S(s),I(s))\Pi_s(dx)ds\\
&\approx \dfrac 1{t}\int_0^t\int_0^1f(x,\varphi(s),0)\Pi_s(dx)ds+\frac 1t\int_0^t \int_0^1xh(x,\varphi(s),0)\Pi_s(dx)ds.
\end{aligned}
\end{equation*}
By the strong law of large numbers for $\varphi(t)$ and $\Pi_t$ from \eqref{Lia}, we obtain that the Lyapunov exponent of $I_{u,v}(t)$  can be
%is
approximated by
\begin{equation}\label{ld}
-c_2+ \int_0^{\infty}\int_{\M} \left(f(x,y,0)\nu(dx)\right)\Phi^*(d\nu)\hat\mu(dy)+\int_0^\infty\int_{\M} \left(\int_0^1 xh(x,y,0)\nu(dx)\right)\Phi^*(d\nu)\hat\mu(dy).
\end{equation}
Since $\mu^*$ is the barycenter of $\Phi^*$, the Lyapunov exponent of $I(t)$ is approximated by
$$
-c_2+ \int_0^{\infty}\int_0^1f(x,y,0)\mu^*(dx)\hat\mu(dy)+\int_0^\infty\int_0^1 xh(x,y,0)\mu^*(dx)\hat\mu(y).
$$
Therefore, we define the threshold $\lambda$ by
\begin{equation}\label{eq-lambda}
\lambda:=-c_2+ \int_0^{\infty}\int_0^1f(x,y,0)\mu^*(dx)\hat\mu(dy)+\int_0^\infty\int_0^1 xh(x,y,0)\mu^*(dx)\hat\mu(dy).
\end{equation}
In the next section, we
%will
prove that the sign of $\lambda$
%will
 characterizes the longtime behavior of the system \eqref{eq-main}.
It is also noted that when $\alpha(t)$ is available,
%observable,
%and
so is \eqref{eq-main-1}, and the permanence or extinction of \eqref{eq-main-1} is also determined by the sign of $\lambda$ defined in \eqref{eq-lambda} (see e.g., \cite{DNY-SPA}).

%Roughly speaking, if $\lambda >0$, whenever $I_{u,v}(t)$ is enough small, $\limsup_{t\to\infty}\frac{\ln I_{u,v}(t)}{t}\approx \lambda>0$ and it leads to $I_{u,v}(t)$ can not be very small in a long time. Conversely, when $\lambda <0$, if the solution starts from a initial point $(u,v)$, where v is sufficiently small then $\limsup_{t\to\infty}\frac{\ln I_{u,v}(t)}{t}\approx \lambda<0$ and which implies $I_{u,v}(t)\to 0$. Therefore, the remaining work is to investigate how the solution enters the region $\{(u,v): v\;\text{is sufficient small}\}$. However, that is just in intuition, the detailed proofs are very technical and complex and need to be carefully done.
%\begin{rem}
%	For $\lambda$ defined as in \eqref{ld}, $\lambda<0$ is equivalent to
%	\begin{equation}\label{r}
%	R:=\dfrac{\int_0^{\infty}f(x,0)f^*(x)dx}{c_2+\int_0^{\infty}\frac 12 g^2(x,0)f^*(x)dx}<1.
%	\end{equation}
%	Therefore, we expect $R$ to be a threshold between the persistence and extinction of \eqref{e2.1} as in the deterministic case.
%	It is noted that since the conditions on $f(s,i)$, it is easy to see that $f(s,i)\leq Fs$ and thus, $$\int_0^{\infty}\Big(f(x,0) - \dfrac 12 g^2(x,0) \Big)f^*(x)dx\leq F\int_0^{\infty}xf^*(x)dx=\dfrac{Fa_1}{b_1}.$$
%	As a result, $\lambda, R$ are well-defined.
%\end{rem}

\section{Characterization of Longtime Properties: Permanence and Extinction}\label{sec:lon}
\subsection{The existence and uniqueness of the solution and preliminary results}
%Let $(\Omega,\F,\{\F_t\}_{t\geq0},\PP)$ be a complete probability space with the filtration $\{\F_t\}_{t\geq0}$ satisfying the usual conditions and $B_1(t),B_2(t),$ $B_3(t)$ be mutually independent Brownian motions.
We begin with the following theorem on the existence and uniqueness of the solution of \eqref{eq-main}
and then proceed with a complete characterization of its positivity and some other important properties.

\begin{thm}\label{thm-exi}
	For any $(u, v,\pi)\in\R_+^2\times\M$, there exists a unique global solution
to
%the
system \eqref{eq-main} with initial value $(u,v,\pi)$.
	The three-component process $\{(S(t), I(t),\Pi_t), t\geq0\}$
	is a Markov process.
\end{thm}

\begin{proof}
	We prove the existence and uniqueness of the solution of \eqref{eq-main} first.
It is noted that although we have assumed $f(s,i)$ is Lipschitz continuous, the coefficient $if(s,i)$ in the system \eqref{eq-main} is not globally Lipschitz in general.
Since the coefficients of the equation are locally Lipschitz continuous, there is a unique solution $(S(t),I(t))$ with the  initial value $(u,v,\pi)\in \R_+^{2}\times\M$, defined on maximal interval  $t\in [0,\tau_e)$, $\tau_e:=\inf\{t\geq 0: S(t)\vee I(t)=\infty\}$ with the convention $\inf\emptyset=\infty$; see e.g., \cite[Theorem 3.8 and Remark 3.10]{MAO06}.
We need to show $\tau_e=\infty$ a.s.
   If we define
$$\tau_k=\inf\Big\{t\geq 0: S(t)\vee I(t)>k\Big\},$$
then $\tau_e=\lim_{k\to\infty}\tau_k$.
Consider $V_1(s,i)=s+i$, then
 %by the generalized It\^o formula,
 we have from definition of $\Lom V_1$
 that
$$
\begin{aligned}
\Lom V_1(s,i,\pi)&=a_1-b_1s-b_2i
%+i\int_0^1h(x,s,i)\pi(dx)
\leq a_1\;\forall (s,i,\pi)\in\R^2_+\times\M.
\end{aligned}
$$
%where $K:=\sup_{[0,1]\times\R}|h(\cdot,\cdot)|<\infty$ as in Assumption \ref{asp-1}.
Hence, by applying It\^o's formula and taking expectation, we obtain
$$
\begin{aligned}
\Ei V_1\big(S(\tau_k\wedge t), I(\tau_k\wedge t)\big)
\leq V_1(u,v)+a_1t,
\end{aligned}
$$
which together with Markov's inequality implies that
$$
\begin{aligned}
\PPi\big\{\tau_k<t\big\}&\leq \PPi\Big\{ V_1\big(S(\tau_k\wedge t), I(\tau_k\wedge t),\alpha(\tau_k\wedge t)\big)\geq k\Big\}
\\&\leq \dfrac{ V_1(u,v)+a_1t}{k}\to 0\;\text{as}\;k\to\infty.
\end{aligned}
$$
Therefore, we have $\PPi\{\tau_e\leq t\}=0$ or $\PPi\{\tau_e>t\}=1$ for all $t>0$. As a consequence, $\PPi\{\tau_e=\infty\}=1$. Hence,  system \eqref{eq-main} has a unique, global, and continuous solution.



We proceed
%continue
 to prove the Markov property. Since $\Pi_t$ is a Markov process and is independent of $B_1(t)$  and $B_2(t)$, the Markov property of the joint process $(S(t),I(t),\Pi_t)$ follows by standard arguments; see
for example,
%this argument can be found in
\cite[Theorem 3.27 and Lemma 3.2]{MAO06} or \cite[Lemma 6.1]{FFK72}. To see why the argument in \cite{MAO06} can be applied,
note that $\Pi_t$ satisfies the stochastic equation \eqref{eq-pi} driven by $\beta(t)$ and the $\sigma$-algebra generated by increments $\{\beta(t_2)-\beta(t_1):t_2\geq t_1\geq t\}$ is independent of $\F_t^y$.
\end{proof}
%\subsection{Preliminary estimates about the expectation}

Next, using Lyapunov functions, we estimate the moments of $S(t)$  and $I(t)$, and obtain some related results. Define $\sigma_*^2:=\max\{\sigma_1^2,\sigma_2^2\}$.

\begin{lem}\label{lem2.1}
The following assertions hold:
\begin{itemize}
\item[{\rm(i)}] For any $0<p<\frac{2\kappa}{\sigma_*^2}$ there is a constant $Q_1$ such that
$$\limsup\limits_{t\to\infty}\Ei (S(t)+ I(t))^{1+p}\leq Q_1\;\; \forall (u,v,\pi) \in \R_+^2\times\M.$$
\item[{\rm(ii)}] For any $\eps>0$, $H>1$, $T>0$, there is $\bar H =\bar H(\eps, H, T)$ such that
$$\PPi \left\{\dfrac{1}{\bar H}  \leq S(t)
 \leq \bar H, \;\;\forall t \in [0,T] \right\}\geq 1-\eps \;\; \text{if}\;\;\; (u,v,\pi) \in [H^{-1},H]\times [0;H]\times\M,$$
and
$$\PPi\{ 0 \leq S(t), I(t) \leq \bar H, \;\;\forall t \in [0,T] \}\geq 1-\eps \;\; \text{if}\;\;\; (u,v,\pi) \in [0,H]\times [0;H]\times\M.$$
\end{itemize}
\end{lem}

\begin{proof}
Consider the Lyapunov function
$V_3(s,i)=(s+i)^{1+p}.$
By directly calculation with
%calculating
the differential operator $\Lom V_3$ and using Assumption \ref{asp-1}, we obtain
\begin{align*}
\Lom V_3(s,i,\pi)=&(1+p)(s+i)^p(a_1-b_1s-b_2i)
%+i\int_0^1xh(x,s,i)\pi(dx))\\
+\frac {p(1+p)}2(s+i)^{p-1}\Big( \sigma_1^2s^2+\sigma_2^2i^2\Big)\\
\leq &-(1+p)(s+i)^{p-1}\Big[ -\frac p2 \sigma_*^2(s+i)^2-a_1(s+i)\big],\;\forall (s,i,\pi)\in\R^2_+\times\M.
\end{align*}
Let
$0<C_4<\frac{p(1+p)}2\sigma_*^2$. By some standard calculations, we get
$$C_5=\sup_{(s,i,\pi)\in \R^2_+\times\M}\left\{\Lom V_3(s,i,\pi)+C_4V_3(s,i)\right\}<\infty.$$
This implies
\begin{equation}\label{10}
\Lom V_3\leq C_5-C_4 V_3.
\end{equation}
Applying \cite[Theorem 5.2, p.157]{MAO} proves part (i) of the lemma. The proof of part (ii) follows
from
%by
part (i) and standard arguments;
%--
see \cite[Lemma 2.1]{DN-JDE}.

%Now, we move to the proof of the part (ii). By \eqref{10}, %we can obtain
% there exist $h_1$,$h_2>0$ (see \cite[Lemma 2.1, p. 45]{DA}) such that for all $(u,v)\in [0,H]\times [0,H]$
%\begin{equation*}
%\mathbb{P}  \Big\{0 \leq S_{u,v} (t) \leq h_1 \;\forall t\in [0,T] \Big\} \geq 1-\frac\varepsilon2\;\text{and}\;\mathbb{P}  \Big\{0 \leq I_{u,v} (t) \leq h_2 \;\forall t\in [0,T] \Big\} \geq 1-\dfrac{\varepsilon}{2}.
%\end{equation*}
%Let
%$$\Omega'_1=\Big\{0 \leq I_{u,v} (t) \leq h_2\;\forall t\in [0,T]\Big\},\;\;\Omega'_1\;\text{may depend on}\;(u,v).$$
% By exponential martingale inequality \cite[Theorem 7.4, p. 44]{MAO} we have $\PP(\Omega'_2)\geq 1-\dfrac{\eps}{2}$, where
%\begin{align*}
%\Omega'_2&=\left\{ -\sigma_1 B_1(t) + \int_0^t\dfrac{I_{u,v}(s)g(S_{u,v}(s),I_{u,v}(s))}{S_{u,v}(s)}dB_3(s)\right.\\
% &\hskip 0.5cm\leq \left.\dfrac{\sigma_1^2 t}{2}+\dfrac12\int_0^t\dfrac{I^2_{u,v}(s)g^2(S_{u,v}(s),I_{u,v}(s))}{S^2_{u,v}(s)}ds + \ln \dfrac{2}{\eps}\;\;\forall t \geq 0\right\},\;\;\Omega'_2\;\text{may depend on}\;(u,v).
%\end{align*}
%Applying  It\^{o}'s formula to the equation \eqref{e2.1} yields that
%\begin{multline}\label{lns}
%%\begin{aligned}
%\ln S_{u,v}(t)=\ln u +\int_0^t \dfrac{a_1}{S_{u,v}(s)}ds -\left(b_1+\dfrac{\sigma_1^2}{2}\right) t-\int_0^t \dfrac{I_{u,v}(s)f(S_{u,v}(s),I_{u,v}(s))}{S_{u,v}(s)}ds
%\\\;\;\;- \dfrac12\int_0^t \dfrac{I^2_{u,v}(s)g^2(S_{u,v}(s),I_{u,v}(s))}{S^2_{u,v}(s)}ds  +\sigma_1B_1(t) - \int_0^t\dfrac{I_{u,v}(s)g(S_{u,v}(s),I_{u,v}(s))}{S_{u,v}(s)}dB_3(s).
%%\end{aligned}
% \end{multline}
%For all $(u,v)\in [H^{-1},H]\times[0,H]$, $\omega\in \Omega'_1\cap\Omega'_2$ and $t \in [0,T]$, by the assumption \ref{conditionfg} we have
%\begin{equation*}
%\begin{aligned}
%\ln S_{u,v}(t)&\geq \ln u -(b_1+\sigma_1^2) t-F\int_0^t I_{u,v}(s)ds -G^2\int_0^t I^2_{u,v}(s)ds-\ln \dfrac 2{\eps}
%\\&\geq \ln H^{-1}-(b_1+\sigma_1^2)T-h_2FT-h_2G^2T-\ln \dfrac 2{\eps}:=\ln h_3^{-1}
%\end{aligned}
%\end{equation*}
%By setting $\bar H=\max\{1, h_1, h_2, h_3\}$ we complete the proof.
\end{proof}


\begin{thm}\label{thm-markov}
	The
%joint
process $(S(t),I(t),\Pi_t)$ is a strong Markov and Feller process.
	Moreover, we have
	$\PPi\{S(t)>0, t>0\}=1$ and
	$\PP_{u,0,\pi}\{I(t)=0, t>0\}=1$,
	$\PPi\{I(t)>0, t>0\}=1$ provided $v>0$.
\end{thm}

\begin{proof}
It is easily seen that
%that
the solution of \eqref{eq-main} is a homogeneous strong Markov and Feller process provided that the coefficients are globally
%globallu
Lipschitz; see e.g., \cite[Theorem 2.9.3]{MAO} and \cite[Section 2.5]{YIN}.
It is noted that the space $\M$ of probability measures in $[0,1]$ endowed with the weak topology can be metricized by the bounded Lipschitz metric defined by
$$
\|\pi_1-\pi_2\|_{\text{BL}}:=\sup\left\{|\pi_1(l)-\pi_2(l)|: \|l\|\leq 1, \sup_{x\neq y\in[0,1]}\frac{|l(x)-l(y)|}{|x-y|}\leq 1\right\}.
$$

Therefore, by using the results in Lemma \ref{lem2.1}, we obtain from the local Lipschitz property of coefficients of \eqref{eq-main} and a truncation
%truncated
argument that $(S(t),I(t),\Pi_t)$ is a homogeneous strong Markov and Feller process.
The details of this truncated argument and this result can be found in \cite[Theorem 5.1]{NYZ17}.



Next, we establish the positivity of solutions. First, suppose that $u,v>0$.
Let us consider  the Lyapunov function $V_2: \R_+^{2}\rightarrow \R_+$
$$V_2(s,i)=\left(s-1-\ln s\right)+(i-1-\ln i).
%\;\;\;\;\;\text{where}\;\;\;\;\;C_1:=\frac {b_2}{F}\cdot
$$
By direct calculations, we have
\begin{equation*}
\begin{aligned}
\Lom V_2(s,i,\pi)&=\left(1-\frac {1}{s}\right)\left(a_1-b_1s-if(x,s,i)-i\int_0^1xh(x,s,i)\pi(dx)\right)+ \frac {\sigma_1^2s^2}{2s^2}
\\
&\;\;\;+\left(1-\frac 1{i}\right)\left(-b_2i+if(x,s,i)+i\int_0^1xh(x,s,i)\pi(dx)\right) +\frac {\sigma_2^2i^2}{2i^2}.
\end{aligned}
\end{equation*}
It follows from Assumption \ref{asp-1} that
$f(x,s,i)=\abs{f(x,s,i)-f(x,0,i)}\leq L_1s$ and
$h(x,s,i)=\abs{h(x,s,i)-h(x,0,i)}\leq L_2s$.
%$$i\int_0^1 h(x,s,i)\pi(dx)\leq b_1s+b_2i.$$
Therefore, it is easily seen that
\begin{equation*}
\begin{aligned}
\Lom V_2(s,i)
&\leq C_1+\frac{f(s,i)i}{s} +\frac{i\int_0^1 h(x,s,i)\pi(dx)}s
\\&\leq C_1+(L_1+L_2)(s+i),
\end{aligned}
\end{equation*}
where $C_1= a_1+ b_1+\frac{\sigma_1^2}2+b_2+\frac{\sigma_2^2}2.$
As a result, if we let $C_2=L_1+L_2+1$ and $C_3=C_1+2C_2\ln C_2+2C_2$ then
\begin{equation}\label{eq-V2}
\begin{aligned}
\Lom V_2(s,i,\pi)-C_2V(s,i)&\leq C_1-s-i+C_2(\ln s+\ln i)+2C_2\\
&\leq C_1+2C_2\ln C_2+2C_2=C_3.
\end{aligned}
\end{equation}
%For any $n\in\N$, we define the following truncated functions
%$$f_0^{(n)}(s,i)=f_0(s\wedge n,i)\;;\;f_{1}^{(n)}(s,x,i)=f_1(s\wedge n,x\wedge n,i)\;;\;f_{2}^{(n)}(s,x,i)=f_2(s\wedge n,x\wedge n,i),$$
%and let $\big(S^{(n)}(t),X^{(n)}(t)\big)$ be the solution of equation \eqref{ww3} and \eqref{eq:tran} with $f_0,f_1,f_2$ replaced by $f_0^{(n)},f_1^{(n)},f_2^{(n)}$, respectively.
For $k>1$, denote
$$\eta=\inf\big\{t\geq 0: S(t)\wedge I(t)\leq 0\big\},$$
$$\eta_k=\inf\Big\{t\geq 0: S(t)\wedge I(t)<\dfrac 1k\Big\}.$$
Then $\eta=\lim_{k\to\infty}\eta_k$.
Therefore, by using the same argument as above, we obtain from \eqref{eq-V2} that
% Consider
%$$\hat V_2^{(n)}(s,x,i)=s-c_1^{(n)}-c_1^{(n)}\ln \dfrac s{c_1^{(n)}}+c_2(x-1-\ln x),$$
%where
%$$c_2=\dfrac 1{\kappa_0}\;\text{and}\;c_1^{(n)}=\frac{c_2\check k_d}{\max_{i\in\M}\sup_{s>0}\frac{f_1^{(n)}(s,x,i)}{s}}.$$
%Let
%$$c_2^{(n)}:=\max_{i\in\M}\sup_{s>0}\frac{f_0^{(n)}(s,i)}{s}.$$
%Since $f_1^{(n)}(s,x,i)$ is global Lipschitz and $f_1^{(n)}(0,x,i)=0$, it is readily seen that
%$\displaystyle\sup_{s>0}\frac{f_1^{(n)}(s,x,i)}{s}<\infty$.
%Similarly, we also obtain  $\sup_{s>0}\frac{f_0^{(n)}(s,i)}{s}<\infty$, so
%the above constants are well-defined.
%
%
%By the generalized It\^o formula, we have
%$$
%\begin{aligned}
%\Lom^{(n)} \hat V_2^{(n)}(s,x,i)=&\Big(1-\dfrac {c_1^{(n)}}s\Big)\left(\dfrac {S_0}\theta-\dfrac{f_0^{(n)}(s,i)}\theta-xf_1^{(n)}(s, x,i)\right)+\dfrac{c_1^{(n)}\sigma_1^2(i)}{2}\\
%&+\Big(c_2-\dfrac{c_2}{x}\Big)\left(xf_2^{(n)}(s, x,i)-x\wdt k_d(i) \right)+\dfrac{c_2\sigma_2^2(i)}{2}
%\\
%\leq& \dfrac{S_0+c_1^{(n)}c_2^{(n)}}{\theta}+c_2\hat k_d+\dfrac{c_1^{(n)}\hat\sigma_1^2+c_2\hat \sigma_2^2}{2}\\
%&+x\Big(c_2f_2^{(n)}(s,x,i)-f_1^{(n)}(s,x,i)\Big)+x\Big(\dfrac{c_1^{(n)}f_1^{(n)}(s,x,i)}{s}-c_2k_d(i)\Big) \\
%\leq& \dfrac{S_0+c_1^{(n)}c_2^{(n)}}{\theta}+c_2\hat k_d+\dfrac{c_1^{(n)}\hat\sigma_1^2+c_2\hat \sigma_2^2}{2}:=K^{(n)} \text{ for } s,x>0,
%\end{aligned}
%,$$
%where the operator $\Lom^{(n)}$ is defined as $\Lom$ with $f_0$, $f_1$, and $f_2$ replaced by $f_0^{(n)}$, $f_1^{(n)}$, and $f_2^{(n)}$, respectively.
%Applying It\^o's formula again, we have
%$$
%\begin{aligned}
%\E_{s,x,i} &\hat V_2^{(n)}\big(S^{(n)}(\eta_k^{(n)}\wedge t), X^{(n)}(\eta_k^{(n)}\wedge t),\alpha(\eta_k^{(n)}\wedge t)\big)\\
%=&\hat V_2^{(n)}(s,x,i)+\E_{s,x,i}\int_0^{\eta_k^{(n)}\wedge t}\Lom \hat V_2^{(n)}\big(S^{(n)}(u), X^{(n)}(u),\alpha(u)\big)du\\
%\leq& \hat V_2^{(n)}(s,x,i)+K^{(n)}t.
%\end{aligned}
%$$
%Since the definition of $\hat V_2^{(n)}$, we observe that if $\eta_k^{(n)}<t$, then
%$$
%\hat V_2^{(n)}\big(S^{(n)}(\eta_k^{(n)}\wedge t), X^{(n)}(\eta_k^{(n)}\wedge t),\alpha(\eta_k^{(n)}\wedge t)\big)\geq \big(c_1^{(n)}\ln kc_1^{(n)}- c_1^{(n)}\big)\wedge \big(c_2\ln k-c_2\big) .
%$$
%Hence, we obtain that
%That is,
%$$\PP_{s,x,i}\big\{S^{(n)}(t),X^{(n)}(t)>0:\forall t>0\big\}=1\text{ for any }n\in\N,$$
%and hence
%$$\PP_{s,x,i}\big\{S^{(n)}(t),X^{(n)}(t)>0:\forall t>0,n\in\N\big\}=1.$$
%Now, for any $t>0$ and
%$$\omega\in \{\tau_e=\infty\}\cap\big\{S^{(n)}(t),X^{(n)}(t)>0:\forall t>0,n\in\N\big\},$$
%there exists $n_0=n_0(\omega,t)$ such that
%$$S(y)(\omega)\vee X(y)(\omega)<n_0\;\forall 0\leq y\leq t.$$
\begin{equation*}
\begin{aligned}
\PPi\{\eta_k<t\}\leq \dfrac{ V_2(u,v)+C_3t}{e^{-C_2t}(\ln k-1)}\to 0\;\text{as}\;k\to \infty.
\end{aligned}
\end{equation*}
As a result, $\PPi\{\eta_\infty=\infty\}=1$. This implies that
\begin{equation}\label{e2-thm2.1}
\PPi\big\{S(t)>0: t>0 \big\}=\PPi\big\{I(t)>0: t>0 \big\}=1\;\forall u,v>0.
\end{equation}
If $u>0,v=0$, the result $\PPi\big\{S(t)>0: t>0 \big\}=1$ can be shown similarly. Moreover, it is obvious that $\PP_{u,0,\pi}\big\{I(t)=0: t>0 \big\}=1$.

We are in a position
%left
to consider the case $u=0$ and $ v\geq 0$ and prove the positivity of $S(t)$.
Let $\eps>0$ be sufficiently small such that
\begin{equation}\label{e3-thm2.1}a_1-b_1\tilde u-\tilde v\sup_{x\in[0,1]}\big(f(x,\tilde u,\tilde v)+h(x,\tilde u,\tilde v)\big)\geq \frac{a_1}{2},
\end{equation}
for any $(\tilde u,\tilde v)\in\R^2$ satisfying
$\tilde u+|\tilde v-v|<\eps$.
Such an $\eps$ exists due to Assumption \ref{asp-1}.
Set $$\tilde\tau_1=\inf\{t>0:S(t)+|I(t)-v|\geq\eps\}.$$
By the continuity of $(S(t), I(t))$ it is clear that $\PP_{0,v,\pi}\{\tilde\tau_1>0\}=1$.
It follows
from \eqref{e3-thm2.1} that
$$a_1-b_1S(t)-I(t)\int_0^1f(x,S(t),I(t))\Pi_t(dx)+I(t)\int_0^1xh(x,S(t),I(t))\Pi_t(dx)>0\;\text{if}\;t\in(0,\tilde\tau_1].$$
This and the variation of constants formula (see \cite[Chapter 3]{MAO})
%we can write $S(t)$ in the form
%\begin{equation}\label{e4-thm2.1}\barray
%\ad
%S(t)=\Phi(t)\left[\int_0^t\Phi^{-1}(u)\left(\dfrac{S_0-f_0\big(S(u),
%	\alpha(u)\big)}\theta-X(u)f_1\big(S(u),X(u),\alpha(u)\big)\right)du\right] \\
%\aad \ \hfill \text{ for } t\in[0,\tilde\tau_1),\earray
%\end{equation}
%where
%$\Phi(t)=\exp\left(-\int_0^t\dfrac{\sigma_1^2(\alpha(u))}2du+\int_0^t\sigma_1(\alpha(u))dW_2(u)\right).$
imply that
$$\PP_{0,v,\pi}\{S(t)>0, t\in(0,\tilde\tau_1]\}=1,$$
which combined with \eqref{e2-thm2.1}
and the strong Markov property of $(S(t),X(t),\Pi(t))$
yields that
\begin{equation*}
\PP_{0,v,\pi}\{S(t)>0, t\in(0,\infty)\}=1.
\end{equation*}
The theorem is therefore proved.
%The solution is also a strong Markov process; see \cite{XM, ZY}.
\end{proof}

\subsection{Extinction}\label{secext}
Consider the case
$\lambda<0$.
We shall show that the number of the infected individuals $I(t)$ tends to zero at an exponential rate while the number of the susceptible individuals $S(t)$ converges to $\varphi(t)$.

\begin{thm}\label{R<1}
Assume that $\lambda<0$. Then for any initial point $(u,v,\pi)\in\mathbb{R}_+^{2,\circ}\times\M$,
the number of the infected individuals $I(t)$ tends to zero  at an exponential rate, i.e., %Further, %, in particular
$$\PPi\left\{\limsup\limits_{t\to\infty}\dfrac{\ln I(t)}{t}= \lambda\right\}=1,$$
and the susceptible class $S(t)$ converges weakly to the solution $\varphi(t)$ on the boundary.
%$$\limsup\limits_{t\to\infty}\dfrac{\ln \abs{S_{u,v}(t)-\varphi_u(t)}}{t}\leq \max\{\lambda,-c_1\}\a.s$$
\end{thm}

In order to prove Theorem \ref{R<1}, we need the following auxiliary results.

\begin{lm}\label{lem2.2}
For any $T,H>1,$ $\eps >0, \theta >0$, there is a $\delta =\delta (H,T,\eps,\theta)$ such that
$$\PPi\{ \tau_{\theta} \geq T \}\geq 1-\eps,\;\;\forall \; (u,v,\pi) \in [0,H]\times (0,\delta]\times\M,$$
where $\tau_{\theta}=\inf\{t\geq 0: I_{u,v}(t)> \theta \}.$
\end{lm}

\begin{proof}
By the exponential martingale inequality \cite[Theorem 7.4, p. 44]{MAO}, we have $\PPi(\Omega_1)\geq 1-\dfrac{\eps}{2}$, where
$$\Omega_1=\left\{ \sigma_2 B_2(t)\leq \dfrac{\sigma_2^2t}{2}+\ln \dfrac{2}{\eps}\;\forall t\geq 0\right\}.$$
%$\Omega''_1$ may depend on $(u,v)$.
In view of part (ii) Lemma \ref{lem2.1}, there exists a $\bar H=\bar H(T,H,\eps)$ such that $\PPi(\Omega_2)\geq 1-\dfrac{\eps}{2}$, where
$$\Omega_2=\{ 0 \leq S(t),I(t) \leq \bar H \;\;\forall t \in [0,T] \}.$$
Applying  It\^{o}'s formula to equation \eqref{eq-main} yields
%implies
that
\begin{equation}\label{lni}
\begin{aligned}
\ln I(t)=&\ln i -c_2t+\int_0^t\int_0^1f(x,S_{u,v}(s),I_{u,v}(s))\Pi_s(dx)ds\\  &+\int_0^t\int_0^1xh(x,S(s),I(s))\Pi_s(dx)ds
+\sigma_2B_2(t).
\end{aligned}
 \end{equation}
Therefore, for any $(s,i,\pi) \in [0,H]\times (0,H]\times\M$ and $\omega \in \Omega_1 \cap \Omega_2$ we have from \eqref{lni} and the Lipschitz continuity of $f$ and $h$ that
$$\ln I(t) < \ln i - b_2T+T\big(2L_1\bar H+2L_2\bar H\big) + \ln \dfrac{2}{\eps},\;\;\forall t\in [0,T]. $$
Hence, we can choose a sufficiently small $\delta =\delta (H,T,\eps,\theta)<H$  such that for all $(s,i,\pi)\in [0,H]\times (0,\delta]\times\M$ and $0\leq t\leq T$, $\ln I(t)<\ln\theta,\forall\omega\in\Omega_1\cap\Omega_2$. The proof is complete.
\end{proof}


\begin{prop}\label{lm3.1}
Suppose that the assumptions from Theorem \ref{R<1} hold. For any $0<\varepsilon<\min\{\frac 15,-\frac{\lambda}{5}\}$ and $H>1$, there exists $\widehat \delta=\widehat \delta(\eps, H)\in (0,H^{-1})$ such that
$$\PPi\left\{\limsup_{t\to\infty}\dfrac{\ln I(t)}{t}=\lambda  \right\}\geq 1-4\eps,\;\forall (u,v,\pi)\in [H^{-1};H]\times(0;\widehat \delta]\times\M_{\Phi^*}.$$
\end{prop}

\begin{proof}
Let $\theta_0=\theta_0(\eps)<\frac{\eps}{L_1+L_2}\wedge \frac {b_1}{L_1+L_2}$ be such that
\begin{equation}\label{eq-theta0}
\frac{a_1+\theta_0(L_1+L_2+L_1\theta_0+L_2\theta_0)}{b_1-(L_1+L_2)\theta_0}-\frac{a_1}{b_1}<\frac{\eps}{L_1+L_2}.
\end{equation}
Consider the following stochastic differential equation
\begin{equation}\label{eq-bphi}
d \bar\varphi(t)=\big(\bar a_1-\bar b_1\bar \varphi(t)\big)dt+\sigma_1 \bar \varphi(t)dB_1(t),
\end{equation}
where $\bar a_1=a_1+L_1+L_2+L_1\theta_0+L_2\theta_0$, $\bar b_1=b_1-(L_1+L_2)\theta_0>0$.
A comparison result shows that $\varphi(t)\leq\bar\varphi(t),t\geq 0\a.s$ provided that $\varphi(0)=\bar \varphi(0)$.
Moreover, the strong law of large numbers yields that
\begin{equation}\label{eq-phibphi}
\begin{aligned}
\lim_{t\to\infty}\frac 1t&\bigg|\int_0^t\int_0^1f(x,\bar\varphi(s),0)\Pi_s(dx)ds+\int_0^t\int_0^1xh(x,\bar \varphi(s),0)\Pi_s(dx)ds\\
&\quad-\int_0^t\int_0^1f(x,\varphi(s),0)\Pi_s(dx)ds-\int_0^t\int_0^1xh(x, \varphi(s),0)\Pi_s(dx)ds\bigg|\\
&\leq \lim_{t\to\infty} \frac1t\int_0^t( L_1+L_2)(\bar \varphi(s)-\varphi(s))ds\\
&=(L_1+L_2)\left(\frac{\bar a_1}{\bar b_1}-\frac{a_1}{b_1}\right)\a.s
\end{aligned}
\end{equation}
Combining \eqref{eq-theta0} and \eqref{eq-phibphi} implies that
\begin{equation}\label{eq-phibphi-1}
\begin{aligned}
\lim_{t\to\infty}\frac 1t&\bigg|\int_0^t\int_0^1f(x,\bar\varphi(s),0)\Pi_s(dx)ds+\int_0^1xh(x,\bar \varphi(s),0)\Pi_s(dx)ds
\\
&\quad-\int_0^t\int_0^1f(x,\varphi(s),0)\Pi_s(dx)ds-\int_0^t\int_0^1xh(x, \varphi(s),0)\Pi_s(dx)ds\bigg| <\eps\a.s
\end{aligned}
\end{equation}
Therefore, there exists $T_1=T_1(\eps)$ such that
$\PP(\Omega_3)\geq 1-\eps$, where
$$
\begin{aligned}
\Omega_3:=&\Bigg\{\bigg|\frac 1t\int_0^t\int_0^1f(x,\bar\varphi(s),0)\Pi_s(dx)ds+\frac1t\int_0^t\int_0^1xh(x,\bar \varphi(s),0)\Pi_s(dx)ds\\
&\quad-\frac 1t\int_0^t\int_0^1f(x,\varphi(s),0)\Pi_s(dx)ds-
\frac1t\int_0^t\int_0^1xh(x, \varphi(s),0)\Pi_s(dx)ds\bigg| <2\eps,\forall t\geq T_1\Bigg\}
\end{aligned}
$$
By definition of $\lambda$ and ergodicity of $\varphi(t), \Pi_t$, we obtain
\begin{equation*}
\PP_{H,\pi}\left\{-c_2+ \lim_{t\to\infty}\dfrac1t\int_0^{t}\int_0^1f(x,\varphi(s),0)\Pi_s(dx)ds+\lim_{t\to\infty}\frac 1t\int_0^t \int_0^1xh(x,\varphi(s),0)\Pi_s(dx)ds= \lambda \right\}=1,
\end{equation*}
where $\PP_{H,\pi}$ indicates the initial values of $(\varphi(t),\Pi_t)$.
As a result, there exists $T_2=T_2(H,\eps)>1$ such that $\PP_{H,\pi}(\Omega_4)\geq 1-\eps$, where
$$\Omega_4=\left\{-c_2+\dfrac1t \int_0^t\int_0^1 f(x,\varphi(s),0)\Pi_s(dx)ds+\frac 1t\int_0^t \int_0^1xh(x,\varphi(s),0)\Pi_s(dx)ds\leq \lambda + \eps,\;\forall t\geq T_2\right\} .$$
In view of the uniqueness
of solutions, we have for all $u\in[0,H]$ that $\varphi_u(s)\leq \varphi_H(s), s \geq 0$ almost surely where the subscript of $\varphi(s)$ indicates the initial value $\varphi(0)$.
This implies that $\PP_{u,\pi}(\Omega_4)\geq 1-\eps$ for all $(u,\pi)\in[0,H]\times\M_{\Phi^*}$.

Since $\lim\limits_{t\to\infty}\dfrac{B_2(t)}{t}=0 \text{ a.s.},$ there is a $T_3=T_3(\eps)>1$ such that $\mathbb{P}(\Omega_5)\geq1-\eps$ where
$$\Omega_5=\left\{\sigma_2\dfrac{B_2(t)}{t}<\eps\;\forall t \geq T_3\right\}.$$
Let $T=\max\{T_1,T_2,T_3\}$. By Lemma \ref{lem2.2}, there is a $\hat\delta=\hat\delta(\eps,H)<\theta_0$ such that for all $(u,v,\pi)\in[0,H]\times (0,\hat\delta]\times\M$,
$
\PPi(\Omega_6)\geq 1-\eps,
$
where
$$
\Omega_6=\{\tau_{\theta_0}\geq T\}.
$$

Now, it follows from \eqref{lni} that $\forall (u,v,\pi) \in [0,H]\times (0,\delta]\times\M_{\Phi^*}$  we have in $\bigcap_{j=3}^{6}\Omega_j$, for all $t\in [T,\tau_{\theta_0}]$ that
\begin{equation}\label{lni2}
\begin{aligned}
\ln I(t)
=&\ln v -c_2t +  \!\!\int_0^{t}\!\Big(\!\int_0^1\! f(x,S(s),I(s))\Pi_s(dx) + \!\! \int_0^1\! xh(x,S(s),I(s))\Pi_s(dx)\!\Big)ds+\sigma_2B_2(t)\\
=&\ln v -c_2t +\int_0^{t}\left(\int_0^1 f(x,S(s),0)\Pi_s(dx) + \int_0^1xh(x,S(s),0)\Pi_s(dx)\right)ds+\sigma_2B_2(t)
\\&+ \int_0^{t}\bigg (\int_0^1f(x,S(s),I(s))\Pi_s(dx) + \int_0^1xh(x,S(s),I(s))\Pi_s(dx)
\\
&\qquad\qquad-\int_0^1f(x,S(s),0)\Pi_s(dx) - \int_0^1xh(x,S(s),0)\Pi_s(dx)\bigg)ds\\
\leq&\ln v -c_2t +  \int_0^{t}\left( \int_0^1f(x,\bar \varphi(s),0)\Pi_s(dx) + \int_0^1xh(x,\bar \varphi(s),0)\Pi_s(dx)\right)ds\\
&\qquad +\eps t+(L_1+L_2)\int_0^t I(s)ds\\
\leq&\ln v -c_2t +  \int_0^{t}\left( \int_0^1f(x,\varphi(s),0)\Pi_s(dx) + \int_0^1xh(x,\varphi(s),0)\Pi_s(dx)\right)ds+3\eps t\\
\leq&\ln v +(\lambda+4\eps)t
<\ln\hat\delta<\ln\theta_0.
\end{aligned}
\end{equation}
In the above, we have used the fact of that whenever $I(t)\leq \theta_0$, one has
$$
\begin{aligned}
I(s)\int_0^1xh(x,S(s),I(s))\Pi_s(dx)&\leq \theta_0(h(0,0,0)+L_2+L_2S(s)+L_2I(s))\\
&\leq \theta_0(L_2+L_2\theta_0)+L_2\theta_0S(s),
\end{aligned}
$$
so $S(s)\leq\bar\varphi(s)$ for all $t\in [0,\tau_{\theta_0}]$.

As a result of \eqref{lni2}, we must have $\tau_{\theta_0}=\infty$ for $\omega\in\bigcap_{j=3}^6\Omega_j$.
We obtain this claim by a contradiction argument as follows. If the claim is false then we have a set $\Omega_7\subset\bigcap_{i=3}^6\Omega_i$
with $\PP(\Omega_7) > 0$ and $\tau_{\theta_0} <\infty$ for any $\omega\in\Omega_7$. We already proved that $\tau_{\theta_0}>T$ for
$\omega\in\bigcap_{i=3}^6\Omega_i$. Moreover, in view of \eqref{lni2}, we have $I(t) \leq \hat\delta<\theta_0$ for any $t\in[T,\tau_{\theta_0}]$. Because
$I(t)$ is continuous almost surely, for almost all $\omega\in\Omega_7$ we have that $\lim_{t\to\tau_{\theta_0}}I(t) = I(\tau_{\theta_0}) <\hat\delta<\theta_0$,
which is a contradiction. So, $\tau_{\theta_0}=\infty$ for $\omega\in\bigcap_{j=1}^4\Omega_j$.
We deduce from $\tau_{\theta_0}=\infty$ and \eqref{lni2} that
%$$Y(t)\leq Y(0)\exp\left\{\frac{\wdt\lambda t}2\right\}, \text{ for any }t\geq T, \omega\in\bigcap_{j=1}^3\Omega_j, Y(0)=y\leq\gamma_1.$$
$\lim_{t\to\infty} I(t)=0$ for almost $\omega\in\bigcap_{j=3}^6\Omega_j$.

Next, because $I(t)\leq \theta_0$ for any $t\geq 0$ for almost all $\omega	\in\bigcap_{j=3}^6\Omega_j$, we have shown that $S(t)\leq \bar\varphi(t), \forall t\geq 0$ almost surely in $\bigcap_{j=3}^6\Omega_j$.
%a comparison argument (\cite[Theorem 6.1.1]{MAO}) implies that $S(t)\leq \hat S(t)$, where
%\begin{equation}\label{se5}
%d\hat S(t)=\left[\sigma-\left(-\frac{\rho\gamma_0}{\eta}+\delta\right) \hat X(t)\right]dt+\sigma_1 \hat X(t)dW_1(t).
%\end{equation} 	
Similar to \eqref{phi}, since $\bar b_1>0$,
the solution to \eqref{eq-bphi} has a unique invariant measure, say $\bar\mu$.
Then we have from the ergodicity of $\bar \varphi(t)$ that for some small $\hat p>0$,
\begin{equation}\label{se6}
\limsup_{t\to\infty}\frac1t\int_0^t S^{1+\hat p}(u)du\leq \lim_{t\to\infty}\frac1t\int_0^t \bar\varphi^{1+\hat p}(u)du=\int x^{1+\hat p}\bar\mu(dx)<\infty \text{ almost surely in }\bigcap_{j=3}^6\Omega_j.
\end{equation} 	
Using \eqref{se6}, the fact that $\lim_{t\to\infty} I(t)=0$, and the compactness of $\M$,
the family of random occupation measures
$$\wdt U^t_{u,v,\pi}(\cdot):=\frac 1t\int_0^t \1_{\{(S(s),I(s),\Pi_s)\in \cdot\}}ds$$ is tight for almost all $\omega\in\bigcap_{j=1}^3\Omega_j$.
From \cite[Lemma 5.6]{HN18},
with probability 1, any weak limit of $\wdt U^t_{u,v,\pi}(\cdot)$ as $t\to\infty$ (if it exists)
is an invariant probability measure,
% of the process $(X(t),Y(t) )$,
which has support on $[0,\infty)\times\{0\}\times\M.$	
Because $\hat\mu\times\bdelta\times\Phi^*$, where $\bdelta$ is the Dirac measure concentrated at $0$, is an
invariant probability measure on $[0,\infty)\times\{0\}\times\M$, the family
$\wdt U^t_{u,v,\pi}(\cdot)$ converges weakly to $\bmu_0\times\bdelta$ almost surely in $ \bigcap_{j=3}^6\Omega_j$ as $t$ tends to $ \infty$.
One has from the weak
convergence
and the uniform integrability
in
%\red{NOTE: Need correct equation here}
\eqref{se6} that
$$
\begin{aligned}
\lim_{t\to\infty}\frac{\ln I(t)}t=&\lim_{t\to\infty}\frac1t\bigg[
-b_2t-\frac{\sigma_2^2t}2-\int_0^t\int_0^1f(x,S(s),I(s))\Pi_s(dx)ds \\
&\qquad\qquad +\int_0^t\int_0^1xh(x,S(s),I(s))\Pi_s(ds)ds\bigg]+\lim_{t\to\infty}\frac{\sigma_2B_2(t)}t\\
=&\lim_{t\to\infty}\!\int_{\R^2_+\times\mathcal P(\M)}\!\!\Big[\!-c_2 +\int_0^1f(x,y,i)\nu(dx)+\int_0^1 xh(x,y,i)\nu(dx)\!\Big]\wdt U^t_{u,v,\pi}(dy,di,d\nu)\\
=&-c_2 +\int_0^\infty \!\int_\M\!\Big[\int_0^1f(x,y,0)\nu(dx)\!\Big]\Phi^*(d\nu)\hat\mu(dy)\\
&+\int_0^\infty\int_{\M} \Big[\int_0^1 xh(x,y,0)\nu(dx)\Big]\Phi^*(d\nu)\hat\mu(dy)\\
=&\lambda<0,
\end{aligned}
$$
for almost every $\omega\in \bigcap_{j=3}^6\Omega_j$, $(u,v ,\pi)\in [0,H]\times (0,\hat\delta]\times\M_{\Phi^*}$.
The proof is complete by noting that $\PP(\bigcap_{j=3}^6\Omega_j)>1-4\eps$.
\iffalse
%Hence, from \eqref{lni2} again, we have $\forall (u,v,\pi)\in [0,H]\times(0,\hat\delta]\times\M_{\Phi^*}$, $\omega\in \cap_{i=3}^{6}\Omega _i$
%\begin{equation*}
%\PPi\left\{\limsup_{t\to\infty}\dfrac{\ln I(t)}{t}\leq \lambda+4\eps  \right\}\geq 1-4\eps.
%\end{equation*}
%The proof is complete.
\fi
\end{proof}


\begin{proof}[Proof of Theorem \ref{R<1}]
	Let $\eps>0$ be arbitrary.
	In view of Proposition \ref{lm3.1},
	the process $(S(t), I(t))$ is transient (see e.g., \cite{WK} for definition) in $\R^{2,\circ}_+$.
	Thus, the process has no invariant probability measure in $\R^{2,\circ}_+$.
	Thus $\hat\mu\times\bdelta\times\Phi^*$
	is the unique invariant probability measure of $(S(t), I(t), \Pi_t)$   in $\R^2_+\times\M_{\Phi^*}$.
	
	Let $H$ be sufficiently large that $\hat\mu((0,H))>1-\eps$.
	%Because of the tightness of
	Thanks to Lemma \ref{lem2.1} part (i) and compactness of $\M$, the process  $(S(t), I(t), \Pi_t)$ is tight. Consequently,
	the occupation measure
	$$U^t_{u,v,\pi}(\cdot):=\dfrac1t\int_0^t\PPi\left\{(S(s),I(s),\Pi_s)\in\cdot\right\}ds$$
	is tight in $\R^2_+\times\M$.
	Since any weak limit of $U^t_{u,v,\pi}$ as $t\to\infty$ must be an invariant probability measure of $(S(t),I(t),\Pi_t)$ (see \cite{HN18}),
	we have that
	$U^t_{u,v,\pi}$ converges weakly to $\hat\mu\times\bdelta\times\Phi^*$ as $t\to\infty$.
	As a result, for any $\delta>0$, there exists a  $\hat T>0$ such that
	$$U^{\hat T}_{u,v,\pi}((0,H)\times(0,\delta)\times\M_{\Phi^*})>1-\eps,$$
	or equivalently,
	$$\dfrac1{\hat T}\int_0^{\hat T}\PPi\{(S(t),I(t),\Pi_t)\in (0,H)\times(0,\delta)\times \M_{\Phi^*}\}dt>1-\eps.$$
As a result, we have
	$$\PPi\{\hat\tau\leq\hat T\}>1-\eps,$$
	where $\hat\tau=\inf\{t\geq 0: (S(t),I(t),\Pi_t)\in (0,H)\times(0,\delta)\times \M_{\Phi^*}\}$.
	Using the strong Markov property and Proposition \ref{lm3.1},
	we have that
	$$
	\PPi\left\{\lim_{t\to\infty}\dfrac{\ln I(t)}t\leq \lambda+4\eps\right\}\geq 1-\eps,
	$$
	% Letting $\eps\to0$ we obtain that
	%$$
	%\PP_{n,p,i}\left\{\lim_{t\to\infty}\dfrac{\ln I(t)}t=\lambda<0, \lim_{t\to\infty}\dfrac{\ln R(t)}t=\max\{-c_3-0.5\sigma^3_2,\lambda\}<0\right\}=1$$
	for any $(u,v,\pi)\in\R^{2,*}_+\times\M$. Therefore, since $\eps > 0$ is arbitrary, the assertion in convergence of $I(t)$ follows.
Once we have the exponentially fast convergence to $0$ of $I(t)$, the convergence of $S(t)$ to $\varphi(t)$ follows from standard arguments; see, for example, \cite{DANG,DN-JDE}.
\end{proof}

\subsection {Permanence}\label{secper}
In this section, we deal with  the case $\lambda>0$ and prove that the system is permanent in the following sense.
\begin{deff}
	We say that system \eqref{eq-main} is permanent (in mean) if for any initial value $(u,v,\pi)\in\R_+^{2,\circ}\times\M$
	$$
	\liminf_{t\to\infty}\frac 1t\int_0^t \E_{u,v,\pi}S(s)ds>0,\quad \liminf_{t\to\infty}\frac 1t\int_0^t \E_{u,v,\pi}I(s)ds>0.
	$$
\end{deff}


\begin{thm}\label{R>1}
Assume that $\lambda>0$. Then for any initial
%point
data
$(u,v,\pi)\in \mathbb{R}_+^{2,\circ}\times\M$,  system \eqref{eq-main} is permanent (in mean).
%, i.e.,  the solution $(S(t),I(t),\Pi_t)$
 %$\wdt\mu(\cdot)$
 %concentrated on $\mathbb{R}_+^{2,\circ}\times\M$.
%Moreover,
%for any $(u,v,\pi)\in\mathbb{R}^{2,\circ}_+\times\M,$
%\begin{equation*}\label{etv}
%\lim\limits_{t\to\infty} \|P(t, (u,v,\pi), \cdot)-\wdt\mu(\cdot)\|=0,
%\end{equation*}
%where $\|\cdot\|$ is the total variation norm, and $P(t, (u, v,\pi), \cdot)$ is the transition probability of $(S(t), I(t),\Pi_t)$.

%\item[{\rm(b)}] The strong large law number holds, i.e, for any $\pi^*$-integrable $ h: \mathbb{R}^{2,\circ}_+\to\mathbb{R}$, we have
%$$\lim\limits_{t\to\infty}\frac1t\int_0^th(S_{u,v}(s), I_{u,v}(s))ds=\int_{\mathbb{R}^{2,\circ}_+} h(x,y)\pi^*(dx, dy) \;a.s.\,\forall \,(u,v)\in\mathbb{R}^{2,\circ}_+.$$
%\end{itemize}
\end{thm}

\begin{proof}
	We first prove that all invariant measures of $(S(t),I(t),\Pi(t)$ concentrate on $\R^{2,\circ}\times\M$.
	We assume by contradiction that there is no invariant measure on $\R^{2,\circ}_+$ of $(S(t),I(t))$.
Therefore, there is no invariant measure on $\R^{2,*}_+$ since the solutions starting in $\R^{2,*}_+$ will enter and remain in $\R^{2,\circ}_+$ due to Theorem \ref{thm-exi}.
	As a result, $\hat\mu\times\bdelta\times\Phi^*$ (the unique invariant on the boundary $\R_+\times\{0\}\times\M$) is the unique invariant probability measure of the process $\{S(t),I(t),\Pi_t\}$ on $\R^2_+\times\M$.
%	For each initial value $(x,y)\in \R^{2,\circ}_+$, consider the  occupation measure
%	$$
%	\Pi_t^{x,y}=\frac1t\E_{x,y}\int_0^t\1_{\{(X(u),Y(u))\in \cdot\}}du.
%	$$
%	Because of \eqref{e1-thm2.1}, $\{\Pi_t^{x,y}, t\geq 1\}$ is a tight family of probability measures on $\R^2_+$.
	% almost surely.
%	Applying \cite[Lemma 3.4]{Hen18}, the tightness implies that any weak limit of $\Pi_t^{x,y}$ is an invariant probability measure of the process $\{X(t),Y(t)\}$.
%	Since $\bmu_0\times\bdelta$ is the unique invariant probability measure, we have that  $\Pi_t^{x,y}$   converges weakly to   $\bmu_0\times\bdelta$ as $t\to\infty$.
	Therefore, by applying \cite[Lemma 3.4]{HN18}, we have
	\begin{equation}\label{limity}
	\begin{aligned}
	\lim_{t\to\infty} \frac1t&\Ei\int_0^t\int_0^1 f(x,S(s),I(s))\Pi_s(dx)ds\\
	&=\int_{\R^2_+\times\M} \int_0^1f(x,y,i)\hat\mu(dy)\bdelta(di)\Phi^*(d\nu)\\
	&=\int_0^\infty\int_0^1 f(x,y,0)\mu^*(dx)\hat\mu(dy).
	\end{aligned}
	\end{equation}
	and
	\begin{equation}\label{limitx}
	\begin{aligned}
	\lim_{t\to\infty} \frac1t&\Ei\int_0^t \int_0^1 xh(x,S(s),I(s))\Pi_s(dx)ds\\
	=&\int_{\R^2_+\times\M} \int_0^1xh(x,y,i)\nu(dx)\hat\mu(dy)\bdelta(di)\Phi^*(d\nu)\\
	&=\int_{\R_+\times\M} \int_0^1xh(x,y,0)\nu(dx)\hat\mu(dy)\Phi^*(d\nu)\\
	&=\int_0^\infty\int_0^1 xh(x,y,0)\mu^*(dx)\hat\mu(dy).
	\end{aligned}
	\end{equation}
	On the other hand,
	$$
	\begin{aligned}
	\Ei\frac{\ln Y(t)}t=&\E_{x,y}\frac{\ln v}t-c_2+\frac{1}{t}\E_{x,y}\bigg(\int_0^t \int_0^1f(x,S(s),I(s))\Pi_s(dx)ds\\
	&\qquad+\int_0^t\int_0^1xh(x,S(s),I(s))\Pi_s(dx)ds\bigg)
	+ \E_{x,y} \frac{\sigma_2W_2(t)}t.
	\end{aligned}
	$$
%	(Note that $\E_{x,y}\frac{\ln Y(t)}t$ exists because
%	all the expectations on the right-hand side exists. In particular $\E \frac{\sigma_2W_2(t)}t$ exists and equals 0).
	As a result,
	we have
	$$
	\begin{aligned}
	\lim_{t\to\infty}&\Ei\frac{\ln I(t)}t\\
	=&
	-c_2+\lim_{t\to\infty}\E_{x,y }\frac1t\left(\int_0^t\int_0^1 f(x,S(s),I(s))\Pi_s(dx)ds+\int_0^t\int_0^1xh(x,S(s),I(s))\Pi_s(dx)ds\right)\\
	=&-c_2+\int_0^\infty \int_0^1f(x,y,0)\mu^*(dx)\hat\mu(dy)+\int_0^\infty\int_0^1 xh(x,y,0)\mu^*(dx)\hat\mu(dy)\\
	=&\lambda>0.
	\end{aligned}
	$$
	This contradicts the fact that
	$$\lim_{t\to\infty}\Ei\frac{\ln I(t)}t\leq \lim_{t\to\infty}\Ei\frac{ I(t)}t=0$$
	because $\ln y\leq y$ while Lemma \ref{lem2.1} implies $\lim_{t\to\infty}\Ei\frac{ I(t)}t=0$.
	As a result, all invariant measures of $(S(t),I(t),\Pi(t))$ concentrate on $\R^{2,\circ}_+\times\M_{\Phi^*}$ (the existence of an invariant measure follows from Lemma \ref{lem2.1}).
	
	By the moment boundedness in Lemma \ref{lem2.1}, we obtain that there exists a sequence $T_k\to\infty$ as $k\to\infty$ such that
	$$
	\liminf_{t\to\infty}\frac 1t\int_0^t\E_{u,v,\pi}S(s)ds=\lim_{k\to\infty}\frac1{T_k}\int_0^{T_k}\E_{u,v,\pi}S(s)ds,
	$$
	$$\liminf_{t\to\infty}\frac 1t\int_0^t\E_{u,v,\pi}I(s)ds=\lim_{k\to\infty}\frac1{T_k}\int_0^{T_k}\E_{u,v,\pi}I(s)ds.
	$$
	Therefore, by applying \cite[Lemma 3.4]{HN18} and the fact that all invariant measures of $(S(t),I(t),\Pi(t))$ concentrate on $\R^{2,\circ}_+\times\M_{\Phi^*}$, we obtain the desired result.
%	The claim \ref{lambda>0} follows easily from a well-known property of non-degenerate diffusi.
\end{proof}
%The main idea to prove this theorem is similar to one in \cite {DANG}. That is to construct a function $V_*:\R_+^{2,\circ}\rightarrow [1,\infty)$ satisfying
%\begin{equation}\label{e4.2}\E V_*(S_{u,v}(t^*),I_{u,v}(t^*))\leq V_*(u,v)-P_1^*V_*^{\gamma}(u,v)+P_2^*\1_{\{u,v\in \K\}} \;\;\;\;\;\forall (u,v)\in\R_+^{2,\circ},
%\end{equation}
%for some petite set $\K$ and some $\gamma\in (0,1)$, $P_1^*,P_2^*>0,t^*>1$ and then apply \cite[Theorem 3.6]{JR}. We also refer the reader to \cite[ pp.106 -124]{MT} for further details on petite sets. Basing on the definition of the value  $\lambda$ in Section \ref{secext}, we will construct $V_*$ as a sum of the Lyapunov function $V_1(u,v)$ defined  in the Lemma \ref{lem2.1} and the function $\ln^- v:=\max\{-\ln v,0\}.$ %\cite[Lemma 2.6]{DANG}As in the Lemma \ref{petite}, we will show that that every compact set is petite. Therefore, the remaining work is to investigate how fast the solution enters a compact set.
% If the solution starts from a initial point $(u,v)$ with  sufficient small  $v$, the functions $\ln^- I_{u,v}(t^*),\ln^- v$ are utilized to dominate the inequality \eqref{e4.2} (see Propositions \ref{prop2.3} and \ref{prop2.4}) while in the remaining region, the Lyapunov functions $V_1(S_{u,v}(t^*),I_{u,v}(t^*)), V_1(u,v)$ play an important role (by using Lemma \ref{lem2.1}).
%Lemmas \ref {lem2.4}, \ref{lem2.5} and \ref{lem2.6} are auxiliary results needed for Propositions \ref{prop2.3} and \ref{prop2.4}.


\subsection{Hidden Markov chain}
%Discrete state space}
\label{sec:dis}
This section is devoted to the case when the signal process $\alpha(t)$ takes values in a  finite set $\{m_1,\dots,m_{n^*}\}\subset[0,1]$ and admits a unique invariant measure $(\mu^*_1,\dots,\mu^*_n)\in\mathcal S_{n^*}$.
In this case,  system \eqref{eq-main} is replaced by \eqref{eq-main-dis-1}.
% and the Assumption \ref{asp-2} is specialized by
%\begin{asp}\label{asp-3}
%	 The generator $(q_{jk})_{\{n^*\times n^*\}}$ of $\alpha(t)$ is irreducible and $g_i - g_j\neq g_l - g_k$ holds if $i\neq l$ or $j\neq k$.
%\end{asp}
%\red{we should split this into three theorems and explain why we dont prove them again in a remark.}
We first have the following well-posedness and other
preliminary results.

\begin{thm} Consider  system \eqref{eq-main-dis-1}.
	% (see \cite{TW11}).
	%\begin{itemize}
	%	\item
		For any $(u, v,e)\in\R_+^2\times \mathcal S_{n^*}$, there exists a unique global solution
		to system \eqref{eq-main-dis-1} with  initial data $(u,v,e)$.
		The three-component process $\{(S(t), I(t),\Pi_t: t\geq0\}$
		is a Markov-Feller process.
		Moreover, we have
		$\PP_{u,v,e}\{S(t)>0, t>0\}=1$ and
		$\PP_{u,0,\pi}\{I(t)=0, t>0\}=1$,
		$\PP_{u,v,e}\{I(t)>0, t>0\}=1$ provided $v>0$.
\end{thm}

%Let $\mathcal S_{n^*}^{0}$ be the support of the invariant measure of Wonham filter.
\begin{proof}
	The proof of this Theorem is as same as that of Theorems \ref{thm-exi} and \ref{thm-markov}, and is thus omitted.
%then we omit here.
\end{proof}

%Moreover,
To proceed,
we classify the persistence and extinction of  system \eqref{eq-main-dis-1} by
the following threshold $\lambda$:
% will be defined as follows.
\begin{equation}\label{eq-lambda-dis}
\lambda:=-c_2+\sum_{k=1}^{n^*}\mu^*_k\int_0^\infty f(m_k,y,0)\hat\mu(dy)+\sum_{k=1}^{n^*} m_k\mu^*_k\int_0^\infty h(m_k,y,0)\hat\mu(dy).
\end{equation}

%[QQ: The next three theorem may be combined as one with three parts??]

\begin{thm}\label{thm-same}
The following results hold.
	\begin{itemize}
		\item [\rm{(i)}]
If $\lambda<0$ then for any initial point $(u,v,e)\in\mathbb{R}_+^{2,\circ}\times\mathcal S_{n^*}$,
the number of the infected individuals $I(t)$ tends to zero  at an exponential rate, that is, %Further, %, in particular
$$\PP_{u,v,e}\left\{\limsup\limits_{t\to\infty}\dfrac{\ln I(t)}{t}=\lambda\right\}=1,$$
and the susceptible class $S(t)$ converges weakly to the solution $\varphi(t)$ on the boundary.

\item [\rm{(ii)}] If $\lambda>0$, then for any initial  point $(u,v,e)\in \mathbb{R}_+^{2,\circ}\times\mathcal S_{n^*}$,
%the
system \eqref{eq-main-dis-1} is permanent (in mean).
%i.e.,  the solution $(S(t),I(t),e(t))$ has a unique invariant probability $\wdt\mu(\cdot)$ concentrated on $\mathbb{R}_+^{2,\circ}\times\mathcal S_{n^*}$.
\item [\rm{(iii)}]
Moreover, the observable system \eqref{eq-main-dis-1} and original
%(unobservable)
system \eqref{eq-main-1} share the same threshold for the persistence and extinction.
	\end{itemize}
\end{thm}

\begin{proof}
	The proofs of part (i) and (ii) in this Theorem are the same as that of Theorems \ref{R<1}, and \ref{R>1} and, therefore are omitted.
%left to the readers.
Part (iii) follows immediately from the formulation of $\lambda$ given in \eqref{eq-lambda-dis} and the threshold for system \eqref{eq-main-1} given in \cite{DNY-SPA}.
\end{proof}

%\begin{thm}
%	
%If $\lambda>0$, then for any initial  point $(u,v,e)\in \mathbb{R}_+^{2,\circ}\times\mathcal S_{n^*}$,
%%the
%system \eqref{eq-main-dis-1} is permanent, i.e.,  the solution $(S(t),I(t),e(t))$ has a unique invariant probability $\wdt\mu(\cdot)$ concentrated on $\mathbb{R}_+^{2,\circ}\times\mathcal S_{n^*}$.
%%\end{itemize}
%\end{thm}
%
%\begin{proof}
%	The proof of this Theorem is
%the
%%as
%same as that of Theorem \ref{R>1} and is omitted
%%we omit
%here.
%\end{proof}

%\begin{thm}\label{thm-same}
%	The observable system \eqref{eq-main-dis-1} and original
%%(unobservable)
%system \eqref{eq-main-1} share the same threshold for the persistence and extinction.
%\end{thm}
%
%\begin{proof}
%	The proof follows immediately from the formulation of $\lambda$ given in \eqref{eq-lambda-dis} and the threshold for system \eqref{eq-main-1} given in \cite{DNY-SPA}.
%\end{proof}

\begin{rem}
Note that in Theorem \ref{thm-same}, we have collected several results. These results may be presented in separate theorems. Given that they are all related to the same hidden process
and that we have carried out an extensive analysis in the last section, it seemed reasonable to collect these results in one theorem.

	In Section \ref{sec:num}, we  present numerical examples for Theorem \ref{thm-same}. The reader can see that although the observable system \eqref{eq-main-dis-1} may not approximate well the original at every point in time, it preserves the longtime properties (permanence or extinction).
\end{rem}

\newcommand{\pre}{{\rm\tiny pre}}
\section{Discussion and Interpretation}\label{sec:disc}
When a pandemic arises, there are usually multiple options one can take in order to  control it. If we apply an extreme policy to control the disease transmission (i.e., we try to reduce the infection rate to be very small or almost $0$), we can certainly control the pandemic.
%We can see this fact
This can be seen
by looking at $\lambda$ defined in the previous section and noting that if $f(S,I)\approx 0$ then $\lambda\approx -c_2<0$.
However, this type of \textit{highly restrictive policy} may hurt the economy.
It is
%an
important
%is
to
%try and
balance public health and the economic issues.
In our context, this is equivalent to ensuring that $\lambda<0$, which ensures the pandemic is under control, but not making $\lambda$ too negative, since in the process of doing this (lockdowns, bankruptcy of businesses, unemployment) the economy could suffer significantly.

\begin{deff}
	We say a proposed threshold $\lambda_1$ is overcautious if it is greater than the exact threshold, that is $\lambda_1>\lambda$.
	We say a proposed threshold $\lambda_1$ is incautious if it is less than the exact threshold, that is $\lambda_1<\lambda$ .
\end{deff}

\begin{rem}
We say	a threshold $\lambda_1>\lambda$ (the exact one) is an overcautious proposed threshold %We say
%It is overcautious because
because if this threshold is implemented,
%proposed,
we tend to
%will try to
apply a policy to make $\lambda_1<0$. But this may not be necessary because the exact threshold is $\lambda<\lambda_1$.
	Conversely, a threshold $\lambda_1<\lambda$ is an incautious one because reducing $\lambda_1$ to be less 0 may not be enough and the pandemic would
%may
still not be controlled.
\end{rem}

%To make our discussion be more analysis,
Let us consider the case when $\alpha(t)$ takes the finitely many values $0=m_1< m_2<\dots<m_{n^*}=1$. We note that our analysis is still true for the general case when $\alpha$ takes values in $[0,1]$.
Recall that we assume $\alpha(t)$ has a unique invariant (discrete) measure $(\mu_1^*,\dots,\mu_{n^*}^*)$.

Since
%we cannot observe
$\alpha(t)$ is not directly available, we have two options.
One option is to use filtering to estimate $\alpha(t)$ and then consider the corresponding system with filtering. It was shown that this method preserves the longtime behavior of the original system.
In particular, from Section \ref{sec:dis}
the exact threshold for this method is
$$
\lambda=-c_2+\sum_{k=1}^{n^*} \mu^*_k\int_0^\infty f(m_k,y,0)\hat\mu(dy)+\sum_{k=1}^{n^*} m_k\mu^*_k\int_0^\infty h(m_k,y,0)\hat\mu(dy).
$$

Another possible option is to estimate some prediction for $\alpha(t)$ and give that value for $\alpha(t)$.
Let $k_0\in\{1,2,\dots,n^*\}$. Assume that we estimate $\alpha(t)=m_{k_0}$ and consider the system
\begin{equation}\label{eq-main-1-alpha=sth}
\begin{cases}
dS(t)=\Big[a_1-b_1S(t)-I(t)f(m_{k_0},S(t),I(t))\\
\hspace{2cm}-\alpha_0I(t)h(m_{k_0},S(t),I(t))\Big]dt + \sigma_1 S(t) dB_1(t),\\[1ex]
dI(t)=\Big[-b_2I(t) + I(t)f(m_{k_0},S(t),I(t))\\
\hspace{2cm}+m_{k_0}I(t)h(m_{k_0},S(t),I(t))\Big]dt + \sigma_2I(t) dB_2(t).\\
\end{cases}
\end{equation}
Using the results from \cite{DNY-SPA}, the threshold for persistence and extinction of \eqref{eq-main-1-alpha=sth} is given by
$$
\lambda_{pre}=-c_2+\int_0^\infty f(m_{k_0},y,0)\hat\mu(dy)+m_{k_0}\int_0^\infty h(m_{k_0},y,0)\hat\mu(dy).
$$
The following
%two
results tell us when $\lambda_{\pre}$ is an incautious or overcautious threshold.

%[QQ: Maybe, the next two propositions can be written as one proposition with two parts?]

\begin{prop}\label{thm-51}
	\begin{itemize}
		\item [\rm{(i)}]
	If
	$$
	\begin{aligned}
	&\mu_{k_0}^*\int_0^\infty (f(m_{k_0},y,0)+m_{k_0}h(m_{k_0},y,0))\hat\mu(dy)\\
	&\hspace{2cm}<\sum_{k=1}^{n^*}\mu^*_k\int_0^\infty f(m_k,y,0)\hat\mu(dy)+\sum_{k=1}^{n^*} m_k\mu^*_k\int_0^\infty h(m_k,y,0)\hat\mu(dy)%\Big)
	\end{aligned}
	$$
	then $\lambda_{\pre}<\lambda$. As a result, $\lambda_{\pre}$ is an incautious threshold.
	\item [\rm{(ii)}]	If
	$$
	\begin{aligned}
	&\mu^*_{k_0}\int_0^\infty (f(m_{k_0},y,0)+m_{k_0}h(m_{k_0},y,0))\hat\mu(dy)\\
	&\hspace{2cm}>\frac 1{1-\mu_{n^*}^*}\Big(\sum_{k=1}^{n^*-1}\mu^*_k\int_0^\infty f(m_k,y,0)\hat\mu(dy)+\sum_{k=1}^{n^*-1} m_k\mu^*_k\int_0^\infty h(m_k,y,0)\hat\mu(dy) \Big)
	\end{aligned}
	$$
	then $\lambda_{\pre}>\lambda$. As a result, $\lambda_{\pre}$ is an overcautious threshold.
		\end{itemize}

\end{prop}

%\begin{prop}\label{thm-52}	
%	If
%	$$
%	\begin{aligned}
%	&\mu^*_{k_0}\int_0^\infty (f(m_{k_0},y,0)+m_{k_0}h(m_{k_0},y,0))\hat\mu(dy)\\
%	&\hspace{2cm}>\frac 1{1-\mu_{n^*}^*}\Big(\sum_{k=1}^{n^*-1}\mu^*_k\int_0^\infty f(m_k,y,0)\hat\mu(dy)+\sum_{k=1}^{n^*-1} m_k\mu^*_k\int_0^\infty h(m_k,y,0)\hat\mu(dy) \Big)
%	\end{aligned}
%	$$
%	then $\lambda_{\pre}>\lambda$. As a result, $\lambda_{\pre}$ is an overcautious threshold.
%\end{prop}

	It is natural to assume that $f(x,y,0)$ is an increasing function w.r.t. $x$ because the higher infection rate in the group of potentially infected individuals would make the disease spread faster.
	Note that this intuition is natural but not
%necessarily
always true because we are examining
%talking about
functions at the boundary, i.e., there is no infected group.

\begin{asp}\label{asp-natural}
	For each fixed $y$, the functions $f(x,y,0)$ and $h(x,y,0)$ are increasing in $x$.
\end{asp}

Under this natural assumption, we will see that assuming that all potentially infected individuals are infected will lead to an overcautious policy. Conversely, it will be incautious if we assume that all individuals with hidden infection status are free of the disease.
We will make this analysis clearer in the following two subsections.

\subsection{The overcautious case: assuming that all individuals in the hidden group are infected}
%We consider equation \eqref{eq-main-1}.
%Recall that we can not observe directly the process $\alpha(t)$.
Suppose we do not use the filtering to consider the observable problem, make the assumption $\alpha(t)=1, t\geq 0$ and consider the system \eqref{eq-main-1-alpha=sth} with $m_{k_0}=m_{n^*}=1$.
Under Assumption \ref{asp-natural}, this is an overcautious prediction.
The following theorem is consistent with this intuition.
%call this case is overcautious pre-prediction.
%In this case, consider the equation
%\begin{equation}\label{eq-main-1-alpha=1}
%\begin{cases}
%dS(t)=\Big[a_1-b_1S(t)-I(t)f(1,S(t),I(t))-I(t)h(1,S(t),I(t))\Big]dt + \sigma_1 S(t) dB_1(t),\\[1ex]
%dI(t)=\Big[-b_2I(t) + I(t)f(1,S(t),I(t))+I(t)h(1,S(t),I(t))\Big]dt + \sigma_2I(t) dB_2(t).\\
%\end{cases}
%\end{equation}
%
%From Section \ref{sec:dis}
%The threshold for this one will be
%$$
%\lambda_{mis}=-c_2+\int_0^\infty f(1,y,0)\hat\mu(dy)+\int_0^\infty h(1,y,0)\hat\mu(dy).
%$$
%$$
%\lambda_{ori}=-c_2+\int_0^\infty f(x,0)\hat\mu(dx)+\sum_{k=1}^{n^*} m_k\mu^*_k\int_0^\infty h(m_k,x,0)\hat\mu(dx).
%$$

%\begin{thm}\label{thm-53}
%	If
%	$$
%	\begin{aligned}
%	&\int_0^\infty (f(1,y,0)+h(1,y,0))\hat\mu(dy)\\
%	&\hspace{2cm}<\frac 1{1-\mu_{n^*}^*}\Big(\sum_{k=1}^{n^*-1}\mu^*_k\int_0^\infty f(m_k,y,0)\hat\mu(dy)+\sum_{k=1}^{n^*-1} m_k\mu^*_k\int_0^\infty h(m_k,y,0)\hat\mu(dy)\Big)
%	\end{aligned}
%	$$
%	then $\lambda_{mis}<\lambda$. As a result, $\lambda_{mis}$ is an incautious threshold.
%	\end{thm}
\begin{prop}\label{thm-53}	
		Under Assumption \ref{asp-natural},
	\begin{equation}\label{asp-lambda-1}
	\begin{aligned}
	&\int_0^\infty (f(1,y,0)+h(1,y,0))\hat\mu(dy)\\
	&\hspace{1cm}>\frac 1{1-\mu_{n^*}^*}\Big(\sum_{k=1}^{n^*-1}\mu^*_k\int_0^\infty f(m_k,y,0)\hat\mu(dy)+\sum_{k=1}^{n^*-1} m_k\mu^*_k\int_0^\infty h(m_k,y,0)\hat\mu(dy)\Big)
	\end{aligned}
	\end{equation}
As a result, $\lambda_{\pre}>\lambda$ and $\lambda_{\pre}$ is an overcautious threshold.
\end{prop}

\begin{proof}
	The estimate \eqref{asp-lambda-1} follows immediately from Assumption \ref{asp-natural}. We omit the details here.
\end{proof}
%\begin{rem}
%	An interesting thing we have seen from Theorem \ref{thm-51} is that an ``overcautious" prediction may lead an incautious result.
%	One may claim that it is impossible due to that it is natural that $f(1,y,0)$ should be always greater than $f(m_k,y,0)$ whenever $m_k<1$. It should be that if the third variable is not $0$. But what the functions here is considered in the case there is no infected people, i.e., the third variable in these function is $0$. Because of this, it is not necessarily satisfied these natural rules.
%\end{rem}

\subsection{The incautious case: assuming that all individuals in the hidden group are not infected}
If we make the assumption that $\alpha(t)=0, t\geq 0$ and consider the system \eqref{eq-main-1-alpha=sth} with $m_{k_0}=m_1=0$, this is an incautious prediction (under Assumption \ref{asp-natural}).
The following result is consistent with this fact.
%In this case, consider the equation
%\begin{equation}\label{eq-main-1-alpha=0}
%\begin{cases}
%dS(t)=\Big[a_1-b_1S(t)-I(t)f(1,S(t),I(t))\Big]dt + \sigma_1 S(t) dB_1(t),\\[1ex]
%dI(t)=\Big[-b_2I(t) + I(t)f(1,S(t),I(t))\Big]dt + \sigma_2I(t) dB_2(t).\\
%\end{cases}
%\end{equation}
%
%From Section \ref{sec:dis}
%The threshold for this one will be
%$$
%\lambda_{mis}=-c_2+\int_0^\infty f(0,y,0)\hat\mu(dy).
%%+\int_0^\infty h(1,y,0)\hat\mu(dy).
%$$
%$$
%\lambda_{ori}=-c_2+\int_0^\infty f(x,0)\hat\mu(dx)+\sum_{k=1}^{n^*} m_k\mu^*_k\int_0^\infty h(m_k,x,0)\hat\mu(dx).
%$$

\begin{prop}\label{thm-54}
	Under Assumption \ref{asp-natural},
	$$
	\begin{aligned}
	&\int_0^\infty (f(0,y,0)+h(0,y,0))\hat\mu(dy)\\
	&\hspace{2cm}<\frac 1{1-\mu_1^*}\Big(\sum_{k=2}^{n^*}\mu^*_k\int_0^\infty f(m_k,y,0)\hat\mu(dy)+\sum_{k=2}^{n^*} m_k\mu^*_k\int_0^\infty h(m_k,y,0)\hat\mu(dy)\Big)
	\end{aligned}
	$$
As a result, $\lambda_{\pre}<\lambda$ and $\lambda_{\pre}$ is an incautious threshold.
\end{prop}

%\begin{thm}\label{thm-52}	
%	If
%	$$
%	\begin{aligned}
%	&\int_0^\infty (f(0,y,0)+h(0,y,0))\hat\mu(dy)\\
%	&\hspace{2cm}>\frac 1{1-\mu_1^*}\Big(\sum_{k=2}^{n^*}\mu^*_k\int_0^\infty f(m_k,y,0)\hat\mu(dy)+\sum_{k=2}^{n^*} m_k\mu^*_k\int_0^\infty h(m_k,y,0)\hat\mu(dy)\Big)
%	\end{aligned}
%	$$
%	then $\lambda_{mis}>\lambda$. As a result, $\lambda_{mis}$ is an overcautious threshold.
%\end{thm}
%\begin{rem}
%\end{rem}


\section{Examples and Simulations}\label{sec:exp}
\subsection{A Simple Example}
In this section, we consider a simple example.
Assume that all the individuals in the hidden class have the same status (susceptible or infected).
In other words, the signal Markov process $\alpha(t)$ has only two possible states, $0$ or $1$, (corresponding to the case that all individuals in hidden class are disease-free or infected, respectively).
%[QQ: mention $0$ and $1$ each represents what, etc.]
Assume that $\alpha(t)$ has the generator
$$
Q=
\begin{bmatrix}
-q_1 &q_1\\
q_2&-q_2
\end{bmatrix}.
$$
We can only observe $\alpha(t)$ via the observation process $y(t)$ given by
$$
dy(t)=g(\alpha(t))dt+dW(t),\quad y(0)=0,
$$
where $g:\{0,1\}\to\R$ and set $g_1=g(0)$, $g_2=g(1)$. Let $e(t)=\E[\1_{\{\alpha(t)=0\}}|\F_t^y]$.

The dynamics of this epidemic system under the (Wonham) filter is described by
\begin{equation}\label{eq-ex-1}
\begin{cases}
dS(t)=\Big[a_1-b_1S(t)-I(t)f(0,S(t),I(t))e(t)-I(t)f(1,S(t),I(t))(1-e(t))\\
\hspace{2cm}-I(t)h(1,S(t),I(t))(1-e(t))\Big]dt + \sigma_1 S(t) dB_1(t),\\[1ex]
dI(t)=\Big[-b_2I(t) + I(t)f(0,S(t),I(t))e(t)+I(t)f(1,S(t),I(t))(1-e(t))\\
\hspace{2cm}+I(t)h(1,S(t),I(t))(1-e(t))\Big]dt + \sigma_2I(t) dB_2(t),\\
de(t)=[q_2-(q_1+q_2)e(t)]dt+(g_1-g_2)e(t)(1-e(t))d\bar W(t).
%\varphi_k(0)=\varphi_k^0,\quad k=1,\dots,n^*.
\end{cases}
\end{equation}
We will also assume that $q_1,q_2>0$ and $g:=g_1-g_2\neq 0$ since the other cases are trivial.
Consider the equation of the third component
\begin{equation}\label{eq-e}
de(t)=[q_2-(q_1+q_2)e(t)]dt+ge(t)(1-e(t))d\bar W(t),\quad e(0)\in [0,1].
\end{equation}
By using the Lyapunov functional method as in Theorem \ref{thm-exi}, we obtain the following result.

\begin{thm}
	For any initial value in $[0,1]$, the equation \eqref{eq-e} has a unique solution and
	$$\PP\{e(t)\in (0,1), \forall t>0\}=1.$$
\end{thm}
Now, by solving the Fokker-Plank equation, equation \eqref{eq-e} has a unique invariant measure supported in $[0,1]$ with density given by
$$
\phi^*(x)=Ce^{-\frac{d_1}{1-x}}(1-x)^{2(d_1-d_2-1)} e^{-\frac{d_2}{x}}x^{2(d_2-d_1-1)},\;x\in(0,1),
$$
where $d_1=\frac{q_1}{g^2}$,
$d_2=\frac{q_2}{g^2}$, and
$C$ is a normalizing constant.

As developed in the main results, the threshold $\lambda$ is defined by
\begin{equation}\label{eq-ex-lambda}
\begin{aligned}
\lambda
%&=-c_2+\lim_{t\to\infty}\dfrac 1{t}\int_0^tf(\varphi(s),0)ds+\lim_{t\to\infty}\frac 1t\int_0^t h(1,\varphi(s),0)(1-e(s))ds\\
%&=-c_2+ \int_0^{\infty}f(x,0)\hat\mu(dx)+\int_0^\infty\int_0^1 h(1,y,0)(1-x)\phi^*(dx)\hat\mu(dy)\\
=&-c_2+ \frac{q_2}{q_1+q_2}\int_0^{\infty}f(0,y,0)\hat\mu(dy)\\
&+\frac{q_1}{q_1+q_2}\int_0^{\infty}f(1,y,0)\hat\mu(dy)+\frac{q_1}{q_1+q_2}\int_0^\infty h(1,y,0)\hat\mu(dy),
\end{aligned}
\end{equation}
where
%$\varphi(s)$ is the solution of the epidemic system on the boundary as in \eqref{phi},
$\hat\mu$ is the invariant measure with the density given by \eqref{hpp}.

\begin{thm} Consider system \eqref{eq-ex-1} and $\lambda$ as in \eqref{eq-ex-lambda}.
	% (see \cite{TW11}).
	\begin{itemize}
%		\item For any $(u, v,e)\in\R_+^2\times \mathcal S_{n^*}$, there exists a unique global solution
%		to the system \eqref{eq-main} initial value $(u,v,e)$.
%		The three-component process $\{(S(t), I(t),\Pi_t, t\geq0\}$
%		is a Markov-Feller process.
%		Moreover, we have
%		$\PP_{u,v,e}\{S(t)>0, t>0\}=1$ and
%		$\PP_{u,0,\pi}\{I(t)=0, t>0\}=1$,
%		$\PP_{u,v,e}\{I(t)>0, t>0\}=1$ provided $v>0$.
		\item
		If $\lambda<0$ then for any initial point $(u,v,e_0)\in\mathbb{R}_+^{2,\circ}\times\mathcal S_{2}$,
		the number of the infected individuals $I(t)$ tends to zero at an exponential rate, i.e., %Further, %, in particular
		$$\PP_{u,v,e_0}\left\{\limsup\limits_{t\to\infty}\dfrac{\ln I(t)}{t}\leq \lambda\right\}=1,$$
		and the susceptible class $S(t)$ converges weakly to the solution $\varphi(t)$ on the boundary.
		\item If $\lambda>0$, then for any initial  point $(u,v,e_0)\in \mathbb{R}_+^{2,\circ}\times\mathcal S_{2}$, all invariant probability measures of
		the solution $(S(t),I(t),e(t))$
		 concentrate on $\mathbb{R}_+^{2,\circ}\times[0,1]$.
		Moreover, the system \eqref{eq-ex-1} is permanent.
	\end{itemize}
\end{thm}


\subsection{Numerical Examples}\label{sec:num}
In this section, we consider a simple example and provide some numerical simulations.
Consider the equation \eqref{eq-ex-1} with $f(x,s,i)=m_1(x)s$,
%$f(1,s,i)=m_2s$,
 %$h(1,s,i)=\frac{m_2s}{1+s+i}$,
 $h(x,s,i)=\frac{m_2(x)s}{1+s+i}$, $g(x)=x$, ($x\in\{0,1\}$), i.e., consider
\begin{equation}\label{eq-ex-2}
\begin{cases}
dS(t)=\Big[a_1-b_1S(t)-m_1(0)S(t)I(t)e(t)-m_1(1)S(t)I(t)(1-e(t))
\\
\hspace{2cm}-\dfrac{m_2(1)(1-e(t))S(t)I(t)}{1+S(t)+I(t)}\Big]dt + \sigma_1 S(t) dB_1(t),\\[2ex]
dI(t)=\Big[-b_2I(t) + m_1(0)S(t)I(t)e(t)+m_1(1)S(t)I(t)(1-e(t))\\
\hspace{2cm}+\dfrac{m_2(1)(1-e(t))S(t)I(t)}{1+S(t)+I(t)}\Big]dt + \sigma_2I(t) dB_2(t),\\
de(t)=[q_2-(q_1+q_2)e(t)]dt+e(t)(1-e(t))d\bar W(t).
%\varphi_k(0)=\varphi_k^0,\quad k=1,\dots,n^*.
\end{cases}
\end{equation}
The above is the corresponding system
%assumes
%is  assuming
%the precise value of $\alpha(t)$ is known (via
with a filtering of hidden Markov chain $\alpha(t)$, whereas
%the observable version of
the following
%(assumed non-observable in this application)
system is one under hidden process $\alpha(t)$
\begin{equation}\label{eq-ex-3}
\begin{cases}
dS(t)=\Big[a_1-b_1S(t)-m_1(\alpha(t))S(t)I(t)\\
\hspace{2cm}-\dfrac{\alpha(t)m_2(\alpha(t))S(t)I(t)}{1+S(t)+I(t)}\Big]dt + \sigma_1 S(t) dB_1(t),\\[2ex]
dI(t)=\Big[-b_2I(t) + m_1(\alpha(t))S(t)I(t)\\
\hspace{2cm}+\dfrac{\alpha(t)m_2(\alpha(t))S(t)I(t)}{1+S(t)+I(t)}\Big]dt + \sigma_2I(t) dB_2(t).\\
%de(t)=[q_2-(q_1+q_2)e(t)]dt+e(t)(1-e(t))d\bar W(t).
%\varphi_k(0)=\varphi_k^0,\quad k=1,\dots,n^*.
\end{cases}
\end{equation}
Here $\alpha(t)$ is the Markov chain taking values in $\{0,1\}$ with the generator $Q$.

Suppose one does not use the filtering to consider the corresponding observable system \eqref{eq-ex-2}, and considers instead the system under some
%fixed
predictions for $\alpha(t)$.
If one uses the incautious prediction and assumes that $\alpha(t)=0$ (i.e., considers all individuals in the hidden group not to be infected), the corresponding system is
\begin{equation}\label{eq-ex-4}
\begin{cases}
dS(t)=\Big[a_1-b_1S(t)-m_1(0)S(t)I(t)\Big]dt + \sigma_1 S(t) dB_1(t),\\[2ex]
dI(t)=\Big[-b_2I(t) + m_1(0)S(t)I(t)\Big]dt + \sigma_2I(t) dB_2(t).\\
%de(t)=[q_2-(q_1+q_2)e(t)]dt+e(t)(1-e(t))d\bar W(t).
%\varphi_k(0)=\varphi_k^0,\quad k=1,\dots,n^*.
\end{cases}
\end{equation}

If one uses the overcautious prediction and assumes that $\alpha(t)=1$ (i.e., considers all individuals in the hidden group to be infected), the corresponding system is
 \begin{equation}\label{eq-ex-5}
 \begin{cases}
 dS(t)=\Big[a_1-b_1S(t)-m_1(1)S(t)I(t)-\dfrac{m_2(1)S(t)I(t)}{1+S(t)+I(t)}\Big]dt + \sigma_1 S(t) dB_1(t),\\[2ex]
 dI(t)=\Big[-b_2I(t) + m_1(1)S(t)I(t)+\dfrac{m_2(1)S(t)I(t)}{1+S(t)+I(t)}\Big]dt + \sigma_2I(t) dB_2(t).\\
 %de(t)=[q_2-(q_1+q_2)e(t)]dt+e(t)(1-e(t))d\bar W(t).
 %\varphi_k(0)=\varphi_k^0,\quad k=1,\dots,n^*.
 \end{cases}
 \end{equation}

\begin{example}
	Consider \eqref{eq-ex-2} with
	$a_1=0.5,
	b_1=1,
	\sigma_1=1,
	b_2=2,
	\sigma_2=0.5,
	m_1(0)=0.1,
	m_1(1)=4,
	m_2=0.1,
	q_1=5,
	q_2=25.
	$
%By computations, we get
Our computation shows that
$\lambda=-1.7252$, the (exact) threshold determining the longtime behavior (persistence and extinction) for both systems \eqref{eq-ex-2} and \eqref{eq-ex-3}.
Similarly, we can compute
$\lambda_0=-2.0750$, the threshold for system \eqref{eq-ex-4} and
$\lambda_1=0.0241$, the threshold for system \eqref{eq-ex-5}.

Note that $\lambda_1>\lambda>\lambda_0$ and as a result $\lambda_1$ is an overcautious threshold. In this case, if we use system \eqref{eq-ex-5}, we may conclude that the disease may not be controlled although it will indeed be controlled. Some unnecessarily restrictive policy might be chosen and this might lead to economic downturns.
Conversely, $\lambda_0$ is an incautious threshold which would lead to overly optimistic expectations.

	As our theoretical results show the number of infected $I(t)$ will tend to $0$ as $t\to\infty$, i.e., the infected group goes extinct.
	We have also shown that  systems \eqref{eq-ex-2} and \eqref{eq-ex-3} have the same longtime behavior.
	We provide the numerical simulations for this example in Figure \ref{fig-1}.
	
	\begin{figure}[h!tb]
		\centering     %%% not \center
		\subfigure{\includegraphics[width=.5\textwidth]{S_t_-ext-new}}\hspace{-1.5em}
		\subfigure{\includegraphics[width=.5\textwidth]{I_t_-ext-new}}
		\subfigure{\includegraphics[width=.5\textwidth]{Filter-ext-new}}
		\caption{Left top subfigure:
%panel:
Sample paths of $S(t)$, in which the red curve is for
%of the
%non-observable
system under the hidden $\alpha(t)$ in \eqref{eq-ex-3}, the blue curve corresponding to the known $\alpha(t)$ in
%of the observable system
\eqref{eq-ex-2}. Right top subfigure: Sample paths of $I(t)$, the red curve is
%of the non-observable
for
system \eqref{eq-ex-3} with hidden state, the blue curve is
%of the observable system
that of \eqref{eq-ex-2} with known $\alpha(t)$. Bottom
subfigure:
%panel:
the sample path of the Markov chains, the red is of the signal process $\alpha(t)$, the blue is its filter $e(t)$.}\label{fig-1}
	\end{figure}
\end{example}

\begin{example}
	Consider \eqref{eq-ex-2} with
	$a_1=10,
	b_1=1,
	\sigma_1=1,
	b_2=3,
	\sigma_2=1,
	m_1(0)=0.1,
	m_1(1)=2,
	m_2(1)=0.1,
	%m_3=0.1;
	q_1=10,
	q_2=1.
	$
Direct computations yield
	$\lambda=14.8522$, the (exact) threshold determining the longtime behavior (persistence and extinction) for both systems \eqref{eq-ex-2} and \eqref{eq-ex-3}.
	Similarly, we can compute
	$\lambda_{0}=-2.5000$, the threshold for system \eqref{eq-ex-4}
	and $\lambda_1=16.5874$, the threshold for system \eqref{eq-ex-5}.
Because $\lambda_1>\lambda>\lambda_0$, we note that $\lambda_1$ is an overcautious threshold. Conversely, $\lambda_0$ is an incautious threshold. [It is readily seen that  the system is actually permanent, i.e., the disease will not be controlled but $\lambda_0$ recommends the disease will be extinct.
Conversely, $\lambda_1$ is an overcautious threshold which would lead to overly pessimistic expectations.]
	
	Applying our theoretical results to this example, we get that $I(t)$ never converges to $0$ and the system is permanent.
		As before, the systems \eqref{eq-ex-2} and \eqref{eq-ex-3} have the same longtime behavior.
	The numerical simulations of this example are provided in Figures \ref{fig-2} and \ref{fig-3}.
	\begin{figure}[h!tb]
		\centering     %%% not \center
		\subfigure{\includegraphics[width=.5\textwidth]{S_t_-per-new1}}\hspace{-1.5em}
		\subfigure{\includegraphics[width=.5\textwidth]{I_t_-per-new1}}
		\subfigure{\includegraphics[width=.5\textwidth]{Filter-per-new1}}
		\caption{Left top subfigure:
%panel:
Sample path of $S(t)$, the red curve is for the
%non-observable
system under noisy observation of $\alpha(t)$, the blue is
%of
the
%observable
system with precise known value of $\alpha(t)$. Right top subfigure:
%panel:
Sample path of $S(t)$, the red curve is for
%the non-observable
system under noisy observation, the blue curve is
%of the observable
for system with precise known value of $\alpha(t)$. Bottom subfigure:
%panel:
the sample path of the Markov chain, the red curve is of the signal process, the blue curve is its filter.}\label{fig-2}
	\end{figure}


\begin{figure}[h!tb]
	\centering     %%% not \center
	\subfigure{\includegraphics[width=.5\textwidth]{2D-hidden}}\hspace{-1.5em}
	\subfigure{\includegraphics[width=.5\textwidth]{3D-hidden}}
	\subfigure{\includegraphics[width=.5\textwidth]{2D-unhidden}}\hspace{-1.5em}
	\subfigure{\includegraphics[width=.5\textwidth]{3D-unhidden}}
	\caption{Left top and right top subfigures:
		%panel:
The density of the invariant probability measure (the marginal one in the space of $(S(t),I(t))$) of \eqref{eq-ex-3} in 2D and 3D settings, respectively.
Bottom subfigure :
	The density of the invariant probability measure (the marginal one in the space of $(S(t),I(t))$) of \eqref{eq-ex-4} in 2D and 3D settings, respectively.}\label{fig-3}
\end{figure}


\end{example}

%\subsection{Discussion}
%We discuss Theorem \ref{R<1} by providing alternative results in some special cases.
%In case of $g(s,i)=0$, we can use the comparison Theorem \cite[Theorem 1.1, p.437]{IW} to obtain that $S_{u,v}(t)\leq \varphi_u(t)\a.s$ Therefore, we can always dominate $S_{u,v}(t)$ by the ergodic process $\varphi_u(t)$ and some estimates can be simplified. %  so we do need any more condition, and we can have an alternative result for this case
%\begin{thm}
%In case of $g(s,i)=0$, the results in Theorem \ref{R<1} hold without the condition $f(s,0)-\dfrac12 g^2(s,0)$ is monotonic.
%\end{thm}
%\begin{proof}
%The proof can be similarly obtained as in \cite[proof of Theorem 2.1]{NHU2}.
%\end{proof}
%If the function $\dfrac {F(s,i)}{s}=\dfrac{if(s,i)}{s}$ is uniformly bounded, for example, the Beddington–DeAngelis incidence rate or the nonlinear functional response with $l=h$, we also obtain the results as in Theorem \ref{R<1}.
%\begin{thm}
%The results in Theorem \ref{R<1} hold if we replace the condition $f(s,0)-\dfrac 12g^2(s,0)$ is monotonic by the condition $\dfrac{F(s,i)}{s}$ is uniformly bounded.
%\end{thm}
%\begin{proof}
%The proof is similar to \cite{NHU1}.
%\end{proof}

%\subsection{The case $R$=1}
\FloatBarrier

\begin{thebibliography}{99}
{\small
\setlength{\baselineskip}{0.10in}
\parskip=0pt

\bibitem{BCL04} P. Baxendale, P. Chigansky, and R. Liptser, Asymptotic stability of
the Wonham filter: Ergodic and nonergodic signals, {\it SIAM J. Control
Optim.}, {\bf 43} (2004), 643--669.

%\bibitem{CAI} Y. Cai, Y. Kang, M. Banerjee, W. Wang, A stochastic SIRS epidemic model with infectious force under intervention strategies, {\it J. Differential Equations}, {\bf 259} (2015), 7463--7502.

%\bibitem{Lou2017}R. Cui, K.Y. Lam,Y. Lou, Dynamics and asymptotic profiles of steady states of an epidemic model in advective environments.,{\it J. Differential Equations}, {\bf 263} (2017),
%% no. 4,
%2343--2373.

%\bibitem{DANG1} N. H. Dang, G. Yin, Coexistence and Exclusion of Stochastic Competitive Lotka-Volterra Models, {\it J. Differential Equations}, {\bf 262} (2017),
%%, no. 3,
%1192--1225.




\bibitem{DANG} N. T. Dieu, D. H. Nguyen, N. H. Du, G. Yin, Classification of Asymptotic Behavior in A Stochastic SIR Model, {\it SIAM J. Appl. Dyn. Syst.}, {\bf 15} (2016),
%no. 2,
1062--1084.
%\bibitem{DDY} N.H. Du, D.H. Nguyen, G. Yin,  Conditions for permanence and ergodicity of certain stochastic predator-prey models, {\it J. Appl. Probab.}, {\bf 53} (2016),
%no.1
%187--202.

\bibitem{NHU1} N. H. Du, N. N. Nhu, Permanence and extinction of certain stochastic SIR models perturbed by a complex type of noises, {\it Appl. Math. Lett.}, {\bf 64} (2017), 223--230.

\bibitem{NHU2}N. H. Du, N. T. Dieu, N. N. Nhu,  Conditions for Permanence and Ergodicity of Certain SIR Epidemic Models, {\it Acta Appl. Math.}, {\bf 160} (2019), 81--99.

%\bibitem{DA} N. H. Du, N. H. Dang, N. T. Dieu, On stability in distribution of stochastic differential delay equations with Markovian switching, {\it Systems Control Lett.}, {\bf 65 } (2014), 43--49.

\bibitem{DN-JDE} N. H. Du, N. N. Nguyen, Permanence and Extinction for the Stochastic SIR Epidemic Model, {\it J. Differential Equations}, {\bf 269} (2020), 9619--9652.


\bibitem{Duncan}
T.E. Duncan,
{\it Probability densities for diffusion processes with applications to nonlinear filtering theory and detection theory}, Ph.D. Diss., Stanford Univ., 1967.

%\bibitem{Duncan}
% T. E. Duncan, Probability density for diffusion processes with
%applications to nonlinear filtering theory, Ph.D. thesis, Standord
%University, 1967.

\bibitem{FFK72} M. Fujisaki,
% Masatoshi;
G. Kallianpur,
H. Kunita,
% Hiroshi.
Stochastic differential equations for the non linear filtering problem, {\it Osaka J. Math.}, {\bf 9} (1972), %no. 1,
19--40.


\bibitem{Han09} R. Handel, The stability of conditional Markov processes and Markov chains in random environments, {\it Ann. Probab.} {\bf 37} (2009), 1876--1925.
% MR-2561436

\bibitem{HN18} A. Hening, D. Nguyen Coexistence and extinction for stochastic Kolmogorov systems,
{\it Ann. Appl. Probab.}, {\bf 28} (2018), 1893--1942.

\bibitem{HUO} L. Huo, J. Jiang, S. Gong, B. He, Dynamical behavior of a rumor transmission model with Holling-type II functional response in emergency event, {\it Phys. A}, {\bf 450} (2016), 228--240.

\bibitem{IW} N. Ikeda, S. Watanabe, {\it Stochastic differential equations and diffusion processes,} Second edition, North-Holland Publishing Co., Amsterdam, (1989).

%\bibitem{JR}S. F. Jarner,  G. O. Roberts,   Polynomial convergence rates of Markov chains, \emph{Ann. Appl. Prob.,} {\bf 12} (2002), 224--247.

%\bibitem{JIANG} D. Q. Jiang, J. J. Yu, C. Y. Ji, and N. Z. Shi, Asymptotic behavior of global positive solution to a stochastic SIR model, {\it Math. Comput. Model.}, {\bf 54} (2011), 221--232.

%\bibitem{Kad1} A. Kaddar, On the dynamics of a delayed SIR epidemic model with a modified saturated
%incidence rate, \emph{Electron. J. Differential Equations}, {\bf 133} (2009), 1--7.


\bibitem{KS69} G. Kallianpur and C. Striebel, Stochastic differential equations occurring in the
estimation of continuous parameter stochastic processes, {\it Theor. Probability Appl.} {\bf 4}
(1969), 597-622.



\bibitem{KER1} W. O. Kermack, A. G. McKendrick, Contributions to the mathematical theory of epidemics I, {\it
Proc. R. Soc. Lond. Ser. A}, {\bf 115} (1927), 700--721.

\bibitem{KER2} W. O. Kermack, A. G. McKendrick, Contributions to the mathematical theory of epidemics II,
{\it Proc. Roy. Soc. Lond. Ser. A}, {\bf 138} (1932), 55--83.

%\bibitem{Kad2} A. Kaddar, Stability analysis in a delayed SIR epidemic model
%with a saturated incidence rate, \emph{Nonlinear Anal. Model. Control}, {\bf 15(3)} (2010), 299--306.

%\bibitem{IK} K. Ichihara,  H. Kunita, A classification of the second order degenerate elliptic operators and its probabilistic characterization, {\it Z. Wahrsch. Verw. Gebiete,} {30} (1974), 235-254. Corrections  { 39}(1977), 81-84.
%\bibitem{JYS} D.Q. Jiang,   J.J. Yu,  C.Y. Ji,  N.Z. Shi, Asymptotic behavior of global positive solution to a stochastic SIR model, \emph{ Math. Comput. Modell.,} {54}(2011), 221-232.

%\bibitem{JJ} C. Ji,   D. Jiang, Threshold behaviour of a stochastic SIR model, \emph{Appl. Math. Model.}, \textbf{38}(2014), no. 21-22, 5067-5079.

%\bibitem{JJYS} C.Y. Ji,   D.Q. Jiang, Q.S. Yang, N.Z. Shi,  Dynamics of a multigroup SIR epidemic model with stochastic perturbation, \emph{Automatica,} \textbf{48}(2012), 121-131.

%\bibitem{JJS}   C.Y. Ji, D.Q. Jiang,  N.Z. Shi,  The behavior of an SIR epidemic model with stochastic perturbation, \emph{Stochastic Anal. Appl.,} {30} (2012), 755-773.

%\bibitem{Ju} V. Jurdjevic, {\it Geometric Control Theory} (Cambridge
%Studies in Advanced Mathematics vol 52)
%Cambridge University Press 1997.

%\bibitem{KM} W.O. Kermack, A. G. McKendrick, Contributions to the mathematical theory of epidemics, (part I), \emph{Proc. R. Soc. Lond. Ser. A,} {115} (1927), 700-721.

%\bibitem{KM1}W. O. Kermack, A. G. McKendrick, Contributions  to  the  mathematical  theory  of
%epidemics,
%(part  II), \emph{ Proc.  Roy.  Sot.  Ser.  A}, {138} (1932), 55-83.

%\bibitem{RK1} R.Z. Khas'minskii, \emph{Stochastic Stability of Differential Equations},
%2nd Ed., Springer-Verlag Berlin Heidelberg, (2012).

%\bibitem{RK} R.Z. Khas'minskii,  Ergodic properties of recurrent diffusion processes and stabilization of the Cauchy problem for parabolic equations. {\it Theory Probab. Appl}. {\bf5} (1960), 179-196.

\bibitem{WK} W. Kliemann,  Recurrence and invariant measures for degenerate diffusions, {\it Ann. Probab.,} {\bf 15} (1987),
% no. 2,
690--707.


\bibitem{Kun71} H.
Kunita,  (1971).
Asymptotic behavior of the nonlinear filtering errors of Markov processes.
J.Multivariate Anal.1 365-393.
%MR0301812


%\bibitem{KRW} D.H. Knipl, G. Rost, J.  Wu,  Epidemic spread and variation of peak times in connected regions due to
%travel-related infections-dynamics of an antigravity-type delay differential model. \emph{SIAM J. Appl. Dyn. Syst.}, 12 (4) (2013), 1722-1762.

%\bibitem{aK}  A. Korobeinikov, Lyapunov functions and global stability for SIR and SIRS epidemiological models with non-linear transmission, \emph{Bull. Math. Biol.,} \textbf{68}(2006), 615-626.

%\bibitem{Kor} I. Kortchemski,  A predator-prey SIR type dynamics on large complete graphs with three phase transitions, \emph{Stochastic Process. Appl.,} 125 (2015), no. 3, 886-917.
%\bibitem{LJX} Y. Lin,    D. Jiang, P. Xia, Long-time behavior of a stochastic SIR model, \emph{Appl. Math. Comput.,} {236} (2014), 1-9.

%\bibitem{LAH} A. Lahrouz, L. Omari, D. Kiouach, Global analysis of a deterministic and stochastic nonlinear SIRS epidemic model, {\it Nonlinear Anal. Model. Control}, {\bf 16} (2011),
%%no. 1,
%59--76.

%\bibitem{Dli-17} D. Li, S. Liu, J. Cui, Threshold dynamics and ergodicity of an SIRS epidemic model with Markovian switching., {\it J. Differential Equations}, {\bf 263} (2017),
%%no. 12,
%8873--8915.

%\bibitem{Dli-18} D. Li, S. Liu, J. Cui, Threshold dynamics and ergodicity of an SIRS epidemic model with semi-Markov switching, {\it J. Differential Equations}, {\bf 266} (2019), 3973-4017.

\bibitem{Kushner64}
  H.J.  Kushner, On the differential equations satisfied by conditional probability densities of Markov processes, with applications,
  {\it J. SIAM Control Ser. A,}, {\bf 2} (1964), 106-119.

\bibitem{LS01} R. S. Liptser, A. N. Shiryaev, {\it Statistics of Random Processes, I: Applications of Mathematics}, (New
York), vol. 5. Springer, Berlin (2001). General theory, Translated from the 1974 Russian original by
A.B. Aries, Stochastic Modelling and Applied Probability



%\bibitem{LIN}Y. Lin, D. Jiang, and P. Xia, Long-time behavior of a stochastic SIR model, {\it Appl. Math. Comput.}, {\bf 236} (2014), 1--9.
%\bibitem{MA} X. Ma, Y. Zhou, H. Cao, Global stability of the endemic equilibrium of a discrete SIR epidemic model. Adv. Difference Equ. 2013, 2013:42.
\bibitem{MAO} X. Mao,  \emph{ Stochastic differential equations and their applications}, Horwood Publishing chichester, 1997.

\bibitem{MAO06} X. Mao and C. Yuan, Stochastic Differential Equations with Markovian Switching, Imperial
College Press, London, 2006.

%\bibitem{MAO1} X. Mao, G. Marion, E. Renshaw, Environmental noise suppresses explosion in population dynamics, {\it Stochastic Process. Appl.}, {\bf 97} (2002), 95--110.

%\bibitem{rM} R.M. May,  \emph{Stability and Complexity in Model Ecosystems}, Princeton University Press, New Jersey, 1973.


\bibitem{Mas66} P. Masani, Wiener’s contributions to generalized harmonic analysis,
prediction theory and filter theory, {\it Bull. Amer. Math. Soc.}, {\bf 72} (1996),
73--125.


%\bibitem{MT} S. P. Meyn, R. L. Tweedie,  {\it Markov Chains and Stochastic Stability,} Springer, London,  1993.

%\bibitem{MT2} S. P. Meyn, R. L. Tweedie, Stability of markovian processes II: continuous-time processes and sampled chains. {\it Adv. in Appl. Probab.}, (1993), 487-517.

%\bibitem{NNY20} D. Nguyen, N. Nguyen, G. Yin, Analysis of a spatially inhomogeneous stochastic partial differential equation epidemic model, to appear in {\it J. Appl. Probab.}
%
%\bibitem{NY19} N. Nguyen, G. Yin, Stochastic partial differential equation SIS epidemic models: modeling and analysis, {\it Commun. Stoch. Anal.}, {\bf 13}, 8 (2019).

 \bibitem{Mortensen}
R.E. Mortensen,
Maximum-likelihood recursive nonlinear filtering,
{\it J. Optim. Theory Appl.}, {\bf 2} (1968), 386-394.


\bibitem{NYZ17} D. Nguyen, G. Yin, Z. Chu, Certain Properties Related to Well Posedness of Switching Diffusions, {\it Stochastic Process. Appl.}, {\bf 127} (2017), 3135--3158.

\bibitem{DNY-SPA} D. Nguyen, N. Nguyen, G. Yin, General nonlinear stochastic systems motivated by chemostat models: Complete characterization of long-time behavior, optimal controls, and applications to wastewater treatment,  {\it Stochastic Process. Appl.}  {\bf 130} (2020), no. 8, 4608--4642.

%\bibitem{EN} E. Nummelin, {\it General Irreducible Markov Chains and Non-negative Operations}, Cambridge  Press, (1984).


\bibitem{TW11} %Takashi
T. Tamura,
% · Yûsuke
Y, Watanabe, Hypoellipticity and ergodicity of the Wonham filter as a diffusion process. (English summary)
{\it Appl. Math. Optim.} {\bf 64} (2011), no. 1, 13–36.

\bibitem{RUAN}
S. Ruan, W. Wang, Dynamical behavior of an epidemic model with a nonlinear incidence rate, {\it J. Differential Equations}, {\bf 188} (2003), 135--163.

%\bibitem{SKO}A. V. Skorokhod, {\it Asymptotic Methods in the Theory of Stochastic Differential Equations}, Translated from the Russian by H. H. McFaden, Transl. Math. Monogr. 78, American Mathematical Society, Providence, RI, 1989.

%\bibitem{SHAN} C. Shan, H. Zhu, Bifurcations and complex dynamics of an SIR model with the impact of the number of hospital beds, {\it J. Differential Equations},  {\bf 257} (2014), 1662--1688.

\bibitem{Sin89} Y. G. Sinaı, Kolmogorov’s work on ergodic theory, {\it Ann. Probab.},
{\bf 17} (1989), 833--839.

%\bibitem{SHAN1} C. Shan, Y. Yi, H. Zhu, Nilpotent singularities and dynamics in an SIR type of compartmental model with hospital resources, {\it J. Differential Equations}, {\bf 260} (2016), 4339--4365.

%\bibitem{SUN} R. Sun, Global stability of the endemic equilibrium of multigroup SIR models with nonlinear incidence. {\it Comput. Math. Appl.}, {\bf 60} (2010),
%% no. 8,
%2286--2291.

%\bibitem{WANG} J. Wang, J. Wang, M. Liu, and Y. Li. Global stability analysis of an SIR epidemic model with demographics and time delay on networks. {\it Phys. A}, {\bf 410} (2014), 268--275.

\bibitem{Wil91} D. Williams, {\it Probability with Martingales}, Cambridge University Press, Cambridge, UK, 1991.

\bibitem{Won65} W.M. Wonham,
Some applications of stochastic differential equations to
optimal nonlinear filtering, {\sl SIAM J. Control}, {\bf 2} (1965),
347--369.
%\bibitem{Won65} W.M. Wonham, Some applications of stochastic differential equations to %optimal nonlinear filtering, {\it J. Soc. Ind.
%Appl. Math. Ser. A Control} {\bf 2} (1965) 347--369.

\bibitem{YANG}
Q. Yang, D. Jiang, N. Shi, C. Ji, The ergodicity and extinction of stochastically perturbed SIR and SEIR epidemic models with saturated incidence, {\it J. Math. Anal. Appl.}, {\bf 388} (2012),
% no. 1,
248--271.

\bibitem{Zakai}
  M. Zakai, On the optimal filtering of diffusion processes,
  {\it Z. Wahrsch. Verw. Gebiete},
  {\bf 11} (1969), 230-243.


\bibitem{YIN}
C. Zhu, G. Yin, On strong Feller, recurrence, and weak stabilization of regime-switching
diffusions, {\it SIAM J. Control Optim.}, {\bf 48} (2009),
% no. 3,
2003--2031.

\bibitem{ZHANG} L. Zhang, Z. C. Wang, X. Q. Zhao, Threshold dynamics of a time periodic reaction–diffusion epidemic model with latent period, {\it J. Differential Equations}, {\bf 258} (2015), 3011--3036.
%\bibitem{PZHANG} F.P. Zhang, Z.Z. Li, F. Zhang, Global stability of an SIR epidemic model with constant infectious period, {\it Appl. Math. Comput.}, {\bf 199} (2008), 285--291.





%\bibitem{SBKS} F. Selley, A.  Besenyei,  I.Z. Kiss, P.L. Simon,  Dynamic control of modern, network-based epidemic models.
%\emph{SIAM J. Appl. Dyn. Syst.,} 14 (2015), no. 1, 168-187.

%\bibitem{SK} D. Siegel, H. Kunze, Monotonicity properties of the SIS and SIR epidemic models, \emph{J. Math. Anal. Appl.,} 185(1994), no. 1, 65-85.

%\bibitem{WZ} W. Wang,  X. Q. Zhao,  Basic reproduction numbers for reaction-diffusion epidemic models. \emph{SIAM J. Appl. Dyn. Syst.,} 11 (2012), no. 4, 1652-1673.

%\bibitem{QDN} Q. Yang, D. Jiang, N. Shi, C. Ji,  The ergodicity and extinction of stochastically perturbed SIR and SEIR epidemic models with saturated incidence, \emph{J. Math. Anal. Appl.,} \textbf{388}(2012), no. 1, 248-271.

%\bibitem{JDN} J. Yu, D. Jiang, N. Shi, [??Ningzhong??] Global stability of two-group SIR model with random perturbation, \emph{J. Math. Anal. Appl.,} \textbf{360}(2009), no. 1, 235-244.

%\bibitem{ZD}X. Zhong,   F. Deng,  Extinction and Persistent of a Stochastic Multi-group SIR Epidemic Model, \emph{Journal of Control Science and Engineering,} \textbf{1}(2013) 13-22.

%\bibitem{YWS} Y. Zhou,    W. Zhang,  S. Yuan, Survival and stationary distribution of a SIR epidemic model with stochastic perturbations, \emph{Appl. Math. Comput.,} {244} (1) (2014),  118-131.
}
\end{thebibliography}

\end{document}
